% Recovered from the previously completed PDF; no new mathematical content.
% Prose is editable; formula boxes retain original symbol positions.
% Build from repository root: xelatex -no-shell-escape -halt-on-error -output-directory=... print/fudan/NAME.tex
\documentclass[11pt,a4paper]{article}
\usepackage[margin=23mm,headheight=14pt,headsep=8.5mm,footskip=10mm]{geometry}
\usepackage{amsmath,amssymb}
\DeclareMathAlphabet{\mathsf}{OT1}{cmss}{m}{n}
\usepackage[no-math]{fontspec}
\usepackage{xeCJK}
\setmainfont{Latin Modern Roman}
\setCJKmainfont{FandolSong-Regular.otf}[Path=assets/fonts/fandol/,BoldFont=FandolSong-Bold.otf,ItalicFont=FandolSong-Regular.otf,Script=CustomDefault]
\xeCJKsetup{PunctStyle=quanjiao,CheckSingle=true}
\usepackage{setspace}\setstretch{1.12}
\setlength{\parindent}{2em}\setlength{\parskip}{3pt}
\emergencystretch=6em
\hfuzz=9pt
\usepackage{titlesec}
\titleformat{\section}{\centering\normalfont\Large\bfseries}{}{0pt}{}
\titlespacing*{\section}{0pt}{1.8\baselineskip}{1.0\baselineskip}
\usepackage{fancyhdr}\pagestyle{fancy}\fancyhf{}
\fancyhead[L]{\small 数学手记·高等代数}\fancyhead[R]{\small 19S01–19S07}
\fancyfoot[C]{\thepage}\renewcommand{\headrulewidth}{0.4pt}
\usepackage[hidelinks,unicode]{hyperref}
\hypersetup{pdftitle={19S01–19S07：题目、解答与小结},pdfauthor={Vincent Arqain}}
\newcommand{\RB}[4]{\raisebox{0pt}[#2bp][#3bp]{\hbox to #1bp{#4\hss}}}
\newcommand{\RG}[5]{\hbox to 0pt{\kern#1bp\raisebox{#2bp}[0pt][0pt]{\fontsize{#3bp}{#3bp}\csname font#4\endcsname\char#5}\hss}}
\newcommand{\RT}[4]{\hbox to 0pt{\kern#1bp\raisebox{#2bp}[0pt][0pt]{\fontsize{#3bp}{#3bp}\selectfont#4}\hss}}
\newcommand{\RS}[4]{\hbox to 0pt{\kern#1bp\raisebox{#2bp}[0pt][0pt]{\fontsize{#3bp}{#3bp}\selectfont\ensuremath{#4}}\hss}}
\newcommand{\RR}[4]{\hbox to 0pt{\kern#1bp\raisebox{#2bp}[0pt][0pt]{\rule{#3bp}{#4bp}}\hss}}
\newcommand{\fontmi}{\usefont{OML}{cmm}{m}{it}}
\newcommand{\fontsy}{\usefont{OMS}{cmsy}{m}{n}}
\newcommand{\fontex}{\usefont{OMX}{cmex}{m}{n}}
\newcommand{\fontrm}{\usefont{OT1}{cmr}{m}{n}}
\newcommand{\fontbf}{\usefont{OT1}{cmr}{bx}{n}}
\newcommand{\fontam}{\usefont{U}{msa}{m}{n}}
\newcommand{\fontbb}{\usefont{U}{msb}{m}{n}}
\widowpenalty=10000\clubpenalty=10000\raggedbottom
\begin{document}
\begin{center}{\LARGE\bfseries 19S01–19S07\par}\vspace{8pt}复旦大学高等代数习题集·题目与解答\end{center}
% Original page 1.
\section*{2019S01}
设\allowbreak{}\hskip.1em \RB{8.366}{8.313}{0.500}{%
\RG{0.000}{0.002}{10.9091}{mi}{65}%
}\hskip.1em\allowbreak{} 为\allowbreak{}\hskip.1em \RB{6.729}{5.324}{0.618}{%
\RG{0.000}{0.002}{10.9091}{mi}{110}%
}\hskip.1em\allowbreak{} 阶复方阵，满足

\[
\RB{61.502}{11.425}{3.230}{%
\RT{12.425}{4.501}{7.9701}{T}\RG{22.951}{4.501}{7.9701}{mi}{109}\RG{53.667}{4.501}{7.9701}{mi}{107}\RG{0.000}{-0.003}{10.9091}{rm}{40}\RG{4.243}{-0.003}{10.9091}{mi}{65}\RG{18.709}{-0.003}{10.9091}{rm}{41}\RG{33.970}{-0.003}{10.9091}{rm}{61}\RG{45.485}{-0.003}{10.9091}{mi}{65}\RG{58.787}{-0.003}{10.9091}{mi}{59}%
}
\]
其中\allowbreak{}\hskip.1em \RB{20.453}{8.067}{2.609}{%
\RG{0.000}{-0.004}{10.9091}{mi}{109}\RG{9.578}{-0.004}{10.9091}{mi}{59}\RG{14.411}{-0.004}{10.9091}{mi}{107}%
}\hskip.1em\allowbreak{} 是互异的正整数。证明：\allowbreak{}\hskip.1em \RB{8.366}{8.307}{0.504}{%
\RG{0.000}{-0.004}{10.9091}{mi}{65}%
}\hskip.1em\allowbreak{} 的特征值为\allowbreak{}\hskip.1em \RB{5.518}{7.762}{0.744}{%
\RG{0.000}{-0.004}{10.9091}{rm}{48}%
}\hskip.1em\allowbreak{} 或单位根。证明：设\allowbreak{}\hskip.1em \RB{8.366}{8.315}{0.500}{%
\RG{0.000}{0.004}{10.9091}{mi}{65}%
}\hskip.1em\allowbreak{} 的特征值按代数重数列为\allowbreak{}\hskip.1em \RB{47.080}{8.075}{2.601}{%
\RG{0.000}{0.004}{10.9091}{mi}{21}\RG{6.363}{-1.633}{7.9701}{rm}{49}\RG{11.096}{0.004}{10.9091}{mi}{59}\RG{15.929}{0.004}{10.9091}{mi}{58}\RG{20.773}{0.004}{10.9091}{mi}{58}\RG{25.605}{0.004}{10.9091}{mi}{58}\RG{30.449}{0.004}{10.9091}{mi}{59}\RG{35.293}{0.004}{10.9091}{mi}{21}\RG{41.702}{-1.633}{7.9701}{mi}{110}%
}\hskip.1em\allowbreak{} 。矩阵\allowbreak{}\hskip.1em \RB{8.366}{8.315}{0.500}{%
\RG{0.000}{0.004}{10.9091}{mi}{65}%
}\hskip.1em\allowbreak{} 与\allowbreak{}\hskip.1em \RB{8.366}{8.315}{0.500}{%
\RG{0.000}{0.004}{10.9091}{mi}{65}%
}\hskip.1em\allowbreak{} T有相同的特征多项式，故\allowbreak{}\hskip.1em \RB{12.608}{8.686}{3.223}{%
\RG{0.000}{0.004}{10.9091}{rm}{40}\RG{4.242}{0.004}{10.9091}{mi}{65}%
}\hskip.1em\allowbreak{} T\allowbreak{}\hskip.1em \RB{11.971}{8.686}{3.223}{%
\RG{0.000}{0.004}{10.9091}{rm}{41}\RG{4.242}{3.962}{7.9701}{mi}{109}%
}\hskip.1em\allowbreak{}与\allowbreak{}\hskip.1em \RB{15.911}{8.310}{0.501}{%
\RG{0.000}{-0.001}{10.9091}{mi}{65}\RG{8.182}{3.957}{7.9701}{mi}{109}%
}\hskip.1em\allowbreak{} 有相同的特征值（计入代数重数）。由题设\allowbreak{}\hskip.1em \RB{12.608}{8.681}{3.228}{%
\RG{0.000}{-0.001}{10.9091}{rm}{40}\RG{4.242}{-0.001}{10.9091}{mi}{65}%
}\hskip.1em\allowbreak{} T\allowbreak{}\hskip.1em \RB{39.706}{9.988}{3.228}{%
\RG{0.000}{-0.001}{10.9091}{rm}{41}\RG{4.242}{3.957}{7.9701}{mi}{109}\RG{15.261}{-0.001}{10.9091}{rm}{61}\RG{26.776}{-0.001}{10.9091}{mi}{65}\RG{34.958}{3.957}{7.9701}{mi}{107}%
}\hskip.1em\allowbreak{} ，可知两组复数

\allowbreak{}\hskip.1em \RB{52.687}{8.514}{3.282}{%
\RG{0.000}{-0.005}{10.9091}{mi}{21}\RG{6.364}{4.499}{7.9701}{mi}{109}\RG{6.364}{-2.702}{7.9701}{rm}{49}\RG{14.352}{-0.005}{10.9091}{mi}{59}\RG{19.185}{-0.005}{10.9091}{mi}{58}\RG{24.029}{-0.005}{10.9091}{mi}{58}\RG{28.872}{-0.005}{10.9091}{mi}{58}\RG{33.705}{-0.005}{10.9091}{mi}{59}\RG{38.549}{-0.005}{10.9091}{mi}{21}\RG{44.958}{4.499}{7.9701}{mi}{109}\RG{44.958}{-2.702}{7.9701}{mi}{110}%
}\hskip.1em\allowbreak{} 与\allowbreak{}\hskip.1em \RB{47.468}{10.530}{3.282}{%
\RG{0.000}{-0.005}{10.9091}{mi}{21}\RG{6.364}{4.499}{7.9701}{mi}{107}\RG{6.364}{-2.702}{7.9701}{rm}{49}\RG{11.484}{-0.005}{10.9091}{mi}{59}\RG{16.317}{-0.005}{10.9091}{mi}{58}\RG{21.161}{-0.005}{10.9091}{mi}{58}\RG{25.993}{-0.005}{10.9091}{mi}{58}\RG{30.837}{-0.005}{10.9091}{mi}{59}\RG{35.681}{-0.005}{10.9091}{mi}{21}\RG{42.090}{4.499}{7.9701}{mi}{107}\RG{42.090}{-2.702}{7.9701}{mi}{110}%
}\hskip.1em\allowbreak{}

作为多重集合相等。

设其中非零特征值为\allowbreak{}\hskip.1em \RB{46.304}{5.320}{2.858}{%
\RG{0.000}{-0.002}{10.9091}{mi}{22}\RG{6.573}{-1.638}{7.9701}{rm}{49}\RG{11.305}{-0.002}{10.9091}{mi}{59}\RG{16.138}{-0.002}{10.9091}{mi}{58}\RG{20.982}{-0.002}{10.9091}{mi}{58}\RG{25.825}{-0.002}{10.9091}{mi}{58}\RG{30.658}{-0.002}{10.9091}{mi}{59}\RG{35.502}{-0.002}{10.9091}{mi}{22}\RG{42.121}{-1.638}{7.9701}{mi}{114}%
}\hskip.1em\allowbreak{} ，均计入代数重数。由于\allowbreak{}\hskip.1em \RB{40.066}{8.069}{2.607}{%
\RG{0.000}{-0.002}{10.9091}{mi}{109}\RG{9.578}{-0.002}{10.9091}{mi}{59}\RG{14.411}{-0.002}{10.9091}{mi}{107}\RG{23.479}{-0.002}{10.9091}{sy}{21}\RG{34.995}{-0.002}{10.9091}{rm}{49}%
}\hskip.1em\allowbreak{} ，去掉上述多重集合中的零后，存在\allowbreak{}\hskip.1em \RB{45.535}{8.685}{3.224}{%
\RG{0.000}{0.003}{10.9091}{sy}{102}\RG{5.455}{0.003}{10.9091}{rm}{49}\RG{10.909}{0.003}{10.9091}{mi}{59}\RG{15.742}{0.003}{10.9091}{mi}{58}\RG{20.586}{0.003}{10.9091}{mi}{58}\RG{25.429}{0.003}{10.9091}{mi}{58}\RG{30.262}{0.003}{10.9091}{mi}{59}\RG{35.106}{0.003}{10.9091}{mi}{114}\RG{40.376}{0.003}{10.9091}{sy}{103}%
}\hskip.1em\allowbreak{} 的一个置换\allowbreak{}\hskip.1em \RB{6.686}{5.205}{0.617}{%
\RG{0.000}{0.003}{10.9091}{mi}{27}%
}\hskip.1em\allowbreak{} ，使得

\[
\RB{127.201}{10.535}{5.712}{%
\RG{6.573}{4.504}{7.9701}{mi}{109}\RG{35.681}{4.504}{7.9701}{mi}{107}\RG{0.000}{-0.001}{10.9091}{mi}{22}\RG{17.592}{-0.001}{10.9091}{rm}{61}\RG{29.108}{-0.001}{10.9091}{mi}{22}\RG{72.593}{-0.001}{10.9091}{rm}{40}\RG{76.825}{-0.001}{10.9091}{rm}{49}\RG{85.320}{-0.001}{10.9091}{sy}{20}\RG{96.835}{-0.001}{10.9091}{mi}{105}\RG{103.624}{-0.001}{10.9091}{sy}{20}\RG{115.139}{-0.001}{10.9091}{mi}{114}\RG{120.364}{-0.001}{10.9091}{rm}{41}\RG{124.606}{-0.001}{10.9091}{mi}{58}\RG{6.573}{-2.697}{7.9701}{mi}{105}\RG{35.681}{-3.219}{7.9701}{mi}{27}\RG{40.806}{-3.219}{7.9701}{rm}{40}\RG{44.099}{-3.219}{7.9701}{mi}{105}\RG{46.982}{-3.219}{7.9701}{rm}{41}%
}
\]
任取\allowbreak{}\hskip.1em \RB{3.696}{7.707}{0.624}{%
\RG{0.000}{-0.004}{10.9091}{mi}{105}%
}\hskip.1em\allowbreak{} ，设它在\allowbreak{}\hskip.1em \RB{6.686}{5.198}{0.624}{%
\RG{0.000}{-0.004}{10.9091}{mi}{27}%
}\hskip.1em\allowbreak{} 中所属循环的长度为\allowbreak{}\hskip.1em \RB{5.082}{5.318}{0.624}{%
\RG{0.000}{-0.004}{10.9091}{mi}{115}%
}\hskip.1em\allowbreak{} ，则

\[
\RB{217.267}{11.794}{5.854}{%
\RG{14.064}{7.319}{5.9776}{mi}{115}\RG{51.785}{7.319}{5.9776}{mi}{115}\RG{55.279}{7.319}{5.9776}{sy}{0}\RG{61.036}{7.319}{5.9776}{rm}{49}\RG{91.425}{7.319}{5.9776}{rm}{50}\RG{103.067}{7.319}{5.9776}{mi}{115}\RG{106.560}{7.319}{5.9776}{sy}{0}\RG{112.317}{7.319}{5.9776}{rm}{50}\RG{169.979}{7.319}{5.9776}{mi}{115}\RG{210.183}{7.319}{5.9776}{mi}{115}\RG{6.573}{4.506}{7.9701}{mi}{109}\RG{39.673}{4.506}{7.9701}{mi}{107}\RG{44.288}{4.506}{7.9701}{mi}{109}\RG{86.803}{4.506}{7.9701}{mi}{107}\RG{95.576}{4.506}{7.9701}{mi}{109}\RG{165.357}{4.506}{7.9701}{mi}{107}\RG{205.562}{4.506}{7.9701}{mi}{107}\RG{0.000}{0.002}{10.9091}{mi}{22}\RG{21.584}{0.002}{10.9091}{rm}{61}\RG{33.099}{0.002}{10.9091}{mi}{22}\RG{68.715}{0.002}{10.9091}{rm}{61}\RG{80.230}{0.002}{10.9091}{mi}{22}\RG{119.996}{0.002}{10.9091}{rm}{61}\RG{131.511}{0.002}{10.9091}{sy}{1}\RG{136.344}{0.002}{10.9091}{sy}{1}\RG{141.188}{0.002}{10.9091}{sy}{1}\RG{147.269}{0.002}{10.9091}{rm}{61}\RG{158.784}{0.002}{10.9091}{mi}{22}\RG{187.473}{0.002}{10.9091}{rm}{61}\RG{198.988}{0.002}{10.9091}{mi}{22}\RG{214.673}{0.002}{10.9091}{mi}{58}\RG{170.482}{-0.948}{5.9776}{mi}{115}\RG{91.929}{-1.092}{5.9776}{rm}{50}\RG{6.573}{-2.695}{7.9701}{mi}{105}\RG{205.562}{-2.695}{7.9701}{mi}{105}\RG{39.673}{-3.217}{7.9701}{mi}{27}\RG{44.798}{-3.217}{7.9701}{rm}{40}\RG{48.091}{-3.217}{7.9701}{mi}{105}\RG{50.974}{-3.217}{7.9701}{rm}{41}\RG{165.357}{-3.217}{7.9701}{mi}{27}\RG{174.474}{-3.217}{7.9701}{rm}{40}\RG{177.768}{-3.217}{7.9701}{mi}{105}\RG{180.652}{-3.217}{7.9701}{rm}{41}\RG{86.803}{-3.361}{7.9701}{mi}{27}\RG{96.080}{-3.361}{7.9701}{rm}{40}\RG{99.373}{-3.361}{7.9701}{mi}{105}\RG{102.256}{-3.361}{7.9701}{rm}{41}%
}
\]
因为\allowbreak{}\hskip.1em \RB{30.002}{8.309}{2.859}{%
\RG{0.000}{-0.002}{10.9091}{mi}{22}\RG{6.573}{-1.639}{7.9701}{mi}{105}\RG{12.985}{-0.002}{10.9091}{sy}{54}\RG{12.985}{-0.002}{10.9091}{rm}{61}\RG{24.483}{-0.002}{10.9091}{rm}{48}%
}\hskip.1em\allowbreak{} 且\allowbreak{}\hskip.1em \RB{38.559}{8.309}{2.859}{%
\RG{0.000}{-0.002}{10.9091}{mi}{109}\RG{9.578}{3.957}{7.9701}{mi}{115}\RG{17.023}{-0.002}{10.9091}{sy}{54}\RG{17.023}{-0.002}{10.9091}{rm}{61}\RG{28.538}{-0.002}{10.9091}{mi}{107}\RG{34.560}{3.957}{7.9701}{mi}{115}%
}\hskip.1em\allowbreak{} ，所以

\[
\RB{61.054}{6.481}{9.584}{%
\RG{16.417}{2.816}{5.9776}{mi}{115}\RG{31.617}{2.816}{5.9776}{mi}{115}\RG{6.574}{0.003}{7.9701}{sy}{106}\RG{8.926}{0.003}{7.9701}{mi}{109}\RG{20.408}{0.003}{7.9701}{sy}{0}\RG{26.995}{0.003}{7.9701}{mi}{107}\RG{35.609}{0.003}{7.9701}{sy}{106}\RG{0.000}{-5.747}{10.9091}{mi}{22}\RG{41.489}{-5.747}{10.9091}{rm}{61}\RG{52.988}{-5.747}{10.9091}{rm}{49}\RG{58.459}{-5.747}{10.9091}{mi}{58}\RG{6.574}{-9.004}{7.9701}{mi}{105}%
}
\]
因此每个非零特征值都是单位根，结论成立。

小结：把矩阵等式转化为特征值幂次的多重集合相等，再沿置换循环迭代，得到单位根条件。

\[
\RB{278.451}{7.724}{2.296}{%
\RT{0.000}{-0.003}{7.9701}{题}\RT{8.966}{-0.003}{7.9701}{目}\RT{17.933}{-0.003}{7.9701}{来}\RT{26.899}{-0.003}{7.9701}{源}\RT{35.866}{-0.003}{7.9701}{：}\RT{44.832}{-0.003}{7.9701}{谢}\RT{53.799}{-0.003}{7.9701}{启}\RT{62.765}{-0.003}{7.9701}{鸿}\RT{71.732}{-0.003}{7.9701}{《}\RT{80.698}{-0.003}{7.9701}{复}\RT{89.665}{-0.003}{7.9701}{旦}\RT{98.631}{-0.003}{7.9701}{大}\RT{107.598}{-0.003}{7.9701}{学}\RT{116.564}{-0.003}{7.9701}{高}\RT{125.531}{-0.003}{7.9701}{等}\RT{134.497}{-0.003}{7.9701}{代}\RT{143.464}{-0.003}{7.9701}{数}\RT{152.430}{-0.003}{7.9701}{习}\RT{161.397}{-0.003}{7.9701}{题}\RT{170.363}{-0.003}{7.9701}{集}\RT{179.329}{-0.003}{7.9701}{（}\RT{188.296}{-0.003}{7.9701}{2}\RT{192.779}{-0.003}{7.9701}{0}\RT{197.262}{-0.003}{7.9701}{2}\RT{201.745}{-0.003}{7.9701}{3}\RT{209.214}{-0.003}{7.9701}{版}\RT{218.180}{-0.003}{7.9701}{）}\RT{222.672}{-0.003}{7.9701}{》}\RT{227.156}{-0.003}{7.9701}{，}\RT{236.122}{-0.003}{7.9701}{第}\RT{248.066}{-0.003}{7.9701}{7}\RT{252.549}{-0.003}{7.9701}{4}\RT{260.018}{-0.003}{7.9701}{页}\RT{268.985}{-0.003}{7.9701}{。}%
}
\]
% Original page 2.
\section*{2019S02}
设\allowbreak{}\hskip.1em \RB{8.889}{7.953}{0.738}{%
\RG{0.000}{0.002}{10.9091}{mi}{86}%
}\hskip.1em\allowbreak{} 为二维实线性空间，\allowbreak{}\hskip.1em \RB{19.395}{8.073}{2.876}{%
\RG{0.000}{0.002}{10.9091}{mi}{39}\RG{7.135}{0.002}{10.9091}{mi}{59}\RG{11.967}{0.002}{10.9091}{mi}{32}%
}\hskip.1em\allowbreak{} 是\allowbreak{}\hskip.1em \RB{8.889}{7.953}{0.738}{%
\RG{0.000}{0.002}{10.9091}{mi}{86}%
}\hskip.1em\allowbreak{} 上两个非零线性变换，满足

\[
\RB{65.196}{8.069}{2.880}{%
\RG{0.000}{-0.002}{10.9091}{mi}{39}\RG{7.135}{-0.002}{10.9091}{mi}{32}\RG{17.058}{-0.002}{10.9091}{rm}{43}\RG{27.967}{-0.002}{10.9091}{mi}{32}\RG{35.451}{-0.002}{10.9091}{mi}{39}\RG{45.632}{-0.002}{10.9091}{rm}{61}\RG{57.131}{-0.002}{10.9091}{rm}{48}\RG{62.601}{-0.002}{10.9091}{mi}{58}%
}
\]
证明：若\allowbreak{}\hskip.1em \RB{8.889}{7.955}{0.736}{%
\RG{0.000}{0.004}{10.9091}{mi}{86}%
}\hskip.1em\allowbreak{} 只有平凡的\allowbreak{}\hskip.1em \RB{7.253}{5.326}{2.874}{%
\RG{0.000}{0.004}{10.9091}{mi}{39}%
}\hskip.1em\allowbreak{} –不变子空间，则\allowbreak{}\hskip.1em \RB{8.889}{7.955}{0.736}{%
\RG{0.000}{0.004}{10.9091}{mi}{86}%
}\hskip.1em\allowbreak{} 必有非平凡的\allowbreak{}\hskip.1em \RB{7.427}{8.075}{2.732}{%
\RG{0.000}{0.004}{10.9091}{mi}{32}%
}\hskip.1em\allowbreak{} –不变子空间。

证明：若\allowbreak{}\hskip.1em \RB{7.253}{5.324}{2.876}{%
\RG{0.000}{0.002}{10.9091}{mi}{39}%
}\hskip.1em\allowbreak{} 有实特征值，则相应的一维特征子空间就是非平凡的\allowbreak{}\hskip.1em \RB{7.253}{5.324}{2.876}{%
\RG{0.000}{0.002}{10.9091}{mi}{39}%
}\hskip.1em\allowbreak{} –不变子空间。故\allowbreak{}\hskip.1em \RB{7.253}{5.324}{2.876}{%
\RG{0.000}{0.002}{10.9091}{mi}{39}%
}\hskip.1em\allowbreak{} 没有实特征值，特别地，\allowbreak{}\hskip.1em \RB{7.253}{5.318}{2.882}{%
\RG{0.000}{-0.004}{10.9091}{mi}{39}%
}\hskip.1em\allowbreak{} 可逆。

记\allowbreak{}\hskip.1em \RB{15.336}{7.333}{0.616}{%
\RG{0.000}{0.004}{10.9091}{mi}{116}\RG{6.970}{0.004}{10.9091}{rm}{61}%
}\hskip.1em\allowbreak{} tr\allowbreak{}\hskip.1em \RB{7.253}{5.326}{2.874}{%
\RG{0.000}{0.004}{10.9091}{mi}{39}%
}\hskip.1em\allowbreak{} 、\allowbreak{}\hskip.1em \RB{17.075}{8.075}{0.616}{%
\RG{0.000}{0.004}{10.9091}{mi}{100}\RG{8.709}{0.004}{10.9091}{rm}{61}%
}\hskip.1em\allowbreak{} det\allowbreak{}\hskip.1em \RB{7.253}{5.326}{2.874}{%
\RG{0.000}{0.004}{10.9091}{mi}{39}%
}\hskip.1em\allowbreak{} 。由Cayley–Hamilton定理，

\[
\RB{89.991}{10.305}{2.878}{%
\RG{7.136}{4.505}{7.9701}{rm}{50}\RG{0.000}{-0.000}{10.9091}{mi}{39}\RG{14.293}{-0.000}{10.9091}{sy}{0}\RG{25.202}{-0.000}{10.9091}{mi}{116}\RG{29.140}{-0.000}{10.9091}{mi}{39}\RG{38.702}{-0.000}{10.9091}{rm}{43}\RG{49.611}{-0.000}{10.9091}{mi}{100}\RG{55.284}{-0.000}{10.9091}{mi}{73}\RG{70.427}{-0.000}{10.9091}{rm}{61}\RG{81.925}{-0.000}{10.9091}{rm}{48}\RG{87.397}{-0.000}{10.9091}{mi}{58}\RG{60.084}{-1.688}{7.9701}{mi}{86}%
}
\]
由于\allowbreak{}\hskip.1em \RB{52.401}{8.067}{2.882}{%
\RG{0.000}{-0.004}{10.9091}{mi}{39}\RG{7.135}{-0.004}{10.9091}{mi}{32}\RG{17.664}{-0.004}{10.9091}{rm}{61}\RG{29.179}{-0.004}{10.9091}{sy}{0}\RG{37.664}{-0.004}{10.9091}{mi}{32}\RG{45.148}{-0.004}{10.9091}{mi}{39}%
}\hskip.1em\allowbreak{} ，有\allowbreak{}\hskip.1em \RB{52.847}{9.755}{2.882}{%
\RG{0.000}{-0.004}{10.9091}{mi}{39}\RG{7.136}{3.955}{7.9701}{rm}{50}\RG{11.868}{-0.004}{10.9091}{mi}{32}\RG{22.396}{-0.004}{10.9091}{rm}{61}\RG{33.911}{-0.004}{10.9091}{mi}{32}\RG{41.395}{-0.004}{10.9091}{mi}{39}\RG{48.545}{3.955}{7.9701}{rm}{50}%
}\hskip.1em\allowbreak{} 。将上式分别右乘和左乘\allowbreak{}\hskip.1em \RB{7.427}{8.067}{2.740}{%
\RG{0.000}{-0.004}{10.9091}{mi}{32}%
}\hskip.1em\allowbreak{} 后相减，得到

\[
\RB{46.623}{8.073}{2.876}{%
\RG{0.000}{0.002}{10.9091}{rm}{50}\RG{5.454}{0.002}{10.9091}{mi}{116}\RG{9.392}{0.002}{10.9091}{mi}{39}\RG{16.527}{0.002}{10.9091}{mi}{32}\RG{27.058}{0.002}{10.9091}{rm}{61}\RG{38.557}{0.002}{10.9091}{rm}{48}\RG{44.028}{0.002}{10.9091}{mi}{58}%
}
\]
由\allowbreak{}\hskip.1em \RB{7.253}{5.321}{2.879}{%
\RG{0.000}{-0.001}{10.9091}{mi}{39}%
}\hskip.1em\allowbreak{} 可逆及\allowbreak{}\hskip.1em \RB{27.545}{8.310}{2.857}{%
\RG{0.000}{-0.001}{10.9091}{mi}{32}\RG{10.528}{-0.001}{10.9091}{sy}{54}\RG{10.528}{-0.001}{10.9091}{rm}{61}\RG{22.026}{-0.001}{10.9091}{rm}{48}%
}\hskip.1em\allowbreak{} ，得\allowbreak{}\hskip.1em \RB{23.987}{7.765}{0.741}{%
\RG{0.000}{-0.001}{10.9091}{mi}{116}\RG{6.970}{-0.001}{10.9091}{rm}{61}\RG{18.468}{-0.001}{10.9091}{rm}{48}%
}\hskip.1em\allowbreak{} 。又因\allowbreak{}\hskip.1em \RB{7.253}{5.321}{2.879}{%
\RG{0.000}{-0.001}{10.9091}{mi}{39}%
}\hskip.1em\allowbreak{} 的特征多项式\allowbreak{}\hskip.1em \RB{30.429}{9.757}{1.406}{%
\RG{0.000}{-0.001}{10.9091}{mi}{120}\RG{6.235}{3.957}{7.9701}{rm}{50}\RG{13.391}{-0.001}{10.9091}{rm}{43}\RG{24.300}{-0.001}{10.9091}{mi}{100}%
}\hskip.1em\allowbreak{} 无实根，故\allowbreak{}\hskip.1em \RB{25.742}{8.070}{0.937}{%
\RG{0.000}{-0.001}{10.9091}{mi}{100}\RG{8.695}{-0.001}{10.9091}{mi}{62}\RG{20.224}{-0.001}{10.9091}{rm}{48}%
}\hskip.1em\allowbreak{} 。设\allowbreak{}\hskip.1em \RB{34.879}{10.231}{1.678}{%
\RG{0.000}{-0.001}{10.9091}{mi}{115}\RG{8.144}{-0.001}{10.9091}{rm}{61}\RG{19.659}{9.295}{10.9091}{sy}{112}\RG{28.750}{-0.002}{10.9091}{mi}{100}\RR{28.750}{9.514}{5.678}{0.436}%
}\hskip.1em\allowbreak{} ，于是

\[
\RB{59.447}{10.299}{2.883}{%
\RG{7.136}{4.499}{7.9701}{rm}{50}\RG{40.012}{4.499}{7.9701}{rm}{50}\RG{0.000}{-0.005}{10.9091}{mi}{39}\RG{14.899}{-0.005}{10.9091}{rm}{61}\RG{26.414}{-0.005}{10.9091}{sy}{0}\RG{34.899}{-0.005}{10.9091}{mi}{115}\RG{44.745}{-0.005}{10.9091}{mi}{73}\RG{56.852}{-0.005}{10.9091}{mi}{58}\RG{49.540}{-1.694}{7.9701}{mi}{86}%
}
\]
任取\allowbreak{}\hskip.1em \RB{52.035}{8.314}{2.853}{%
\RG{0.000}{0.003}{10.9091}{rm}{48}\RG{8.485}{0.003}{10.9091}{sy}{54}\RG{8.485}{0.003}{10.9091}{rm}{61}\RG{20.000}{0.003}{10.9091}{mi}{101}\RG{25.079}{-1.633}{7.9701}{rm}{49}\RG{32.842}{0.003}{10.9091}{sy}{50}\RG{43.145}{0.003}{10.9091}{mi}{86}%
}\hskip.1em\allowbreak{} ，令\allowbreak{}\hskip.1em \RB{65.634}{9.762}{3.224}{%
\RG{0.000}{0.003}{10.9091}{mi}{101}\RG{5.080}{-1.633}{7.9701}{rm}{50}\RG{12.842}{0.003}{10.9091}{rm}{61}\RG{24.357}{0.003}{10.9091}{mi}{115}\RG{29.471}{3.962}{7.9701}{sy}{0}\RG{36.058}{3.962}{7.9701}{rm}{49}\RG{40.790}{0.003}{10.9091}{mi}{39}\RG{47.926}{0.003}{10.9091}{rm}{40}\RG{52.169}{0.003}{10.9091}{mi}{101}\RG{57.248}{-1.633}{7.9701}{rm}{49}\RG{61.981}{0.003}{10.9091}{rm}{41}%
}\hskip.1em\allowbreak{} 。因\allowbreak{}\hskip.1em \RB{7.253}{5.325}{2.875}{%
\RG{0.000}{0.003}{10.9091}{mi}{39}%
}\hskip.1em\allowbreak{} 无实特征值，\allowbreak{}\hskip.1em \RB{24.042}{5.325}{2.602}{%
\RG{0.000}{0.003}{10.9091}{mi}{101}\RG{5.079}{-1.633}{7.9701}{rm}{49}\RG{9.811}{0.003}{10.9091}{mi}{59}\RG{14.644}{0.003}{10.9091}{mi}{101}\RG{19.740}{-1.633}{7.9701}{rm}{50}%
}\hskip.1em\allowbreak{} 线性无关。于是\allowbreak{}\hskip.1em \RB{7.253}{5.325}{2.875}{%
\RG{0.000}{0.003}{10.9091}{mi}{39}%
}\hskip.1em\allowbreak{} 在基\allowbreak{}\hskip.1em \RB{35.084}{8.685}{3.224}{%
\RG{0.000}{0.003}{10.9091}{sy}{102}\RG{5.454}{0.003}{10.9091}{mi}{101}\RG{10.534}{-1.633}{7.9701}{rm}{49}\RG{15.266}{0.003}{10.9091}{mi}{59}\RG{20.099}{0.003}{10.9091}{mi}{101}\RG{25.194}{-1.633}{7.9701}{rm}{50}\RG{29.926}{0.003}{10.9091}{sy}{103}%
}\hskip.1em\allowbreak{}下的表示矩阵为

\[
\RB{81.050}{19.587}{14.130}{%
\RG{30.008}{18.650}{10.9092}{ex}{32}\RG{68.002}{18.650}{10.9092}{ex}{33}\RG{38.644}{7.275}{10.9091}{rm}{48}\RG{54.061}{7.275}{10.9091}{sy}{0}\RG{62.547}{7.275}{10.9091}{rm}{49}\RG{0.000}{-0.004}{10.9091}{mi}{70}\RG{11.561}{-0.004}{10.9091}{rm}{61}\RG{23.076}{-0.004}{10.9091}{mi}{115}\RG{78.456}{-0.004}{10.9091}{mi}{58}\RG{38.644}{-7.900}{10.9091}{rm}{49}\RG{58.303}{-7.900}{10.9091}{rm}{48}%
}
\]
\noindent \RB{472.935}{19.588}{58.744}{%
\RG{182.020}{18.651}{10.9092}{ex}{32}\RG{214.402}{18.651}{10.9092}{ex}{33}\RG{191.586}{7.276}{10.9091}{mi}{117}\RG{208.713}{7.276}{10.9091}{mi}{118}\RT{0.000}{-0.004}{10.9091}{设}\RG{14.102}{-0.004}{10.9091}{mi}{32}\RT{24.793}{-0.004}{10.9091}{在}\RT{35.702}{-0.004}{10.9091}{同}\RT{46.612}{-0.004}{10.9091}{一}\RT{57.521}{-0.004}{10.9091}{组}\RT{68.430}{-0.004}{10.9091}{基}\RT{79.339}{-0.004}{10.9091}{下}\RT{90.248}{-0.004}{10.9091}{的}\RT{101.158}{-0.004}{10.9091}{表}\RT{112.067}{-0.004}{10.9091}{示}\RT{122.976}{-0.004}{10.9091}{矩}\RT{133.885}{-0.004}{10.9091}{阵}\RT{144.794}{-0.004}{10.9091}{为}\RG{158.897}{-0.004}{10.9091}{mi}{71}\RG{170.504}{-0.004}{10.9091}{rm}{61}\RT{223.038}{-0.004}{10.9091}{。}\RT{231.972}{-0.004}{10.9091}{由}\RG{246.071}{-0.004}{10.9091}{mi}{70}\RG{254.591}{-0.004}{10.9091}{mi}{71}\RG{264.723}{-0.004}{10.9091}{rm}{43}\RG{274.752}{-0.004}{10.9091}{mi}{71}\RG{283.327}{-0.004}{10.9091}{mi}{70}\RG{294.890}{-0.004}{10.9091}{rm}{61}\RG{306.389}{-0.004}{10.9091}{rm}{48}\RT{311.861}{-0.004}{10.9091}{，}\RT{320.795}{-0.004}{10.9091}{逐}\RT{331.704}{-0.004}{10.9091}{项}\RT{342.614}{-0.004}{10.9091}{比}\RT{353.523}{-0.004}{10.9091}{较}\RT{364.432}{-0.004}{10.9091}{得}\RT{375.341}{-0.004}{10.9091}{到}\RG{389.438}{-0.004}{10.9091}{mi}{119}\RG{400.572}{-0.004}{10.9091}{rm}{61}\RG{412.087}{-0.004}{10.9091}{mi}{118}\RT{417.767}{-0.004}{10.9091}{、}\RG{426.698}{-0.004}{10.9091}{mi}{122}\RG{435.281}{-0.004}{10.9091}{rm}{61}\RG{446.796}{-0.004}{10.9091}{sy}{0}\RG{455.281}{-0.004}{10.9091}{mi}{117}\RT{461.526}{-0.004}{10.9091}{，}\RG{190.657}{-7.899}{10.9091}{mi}{119}\RG{208.766}{-7.899}{10.9091}{mi}{122}\RT{0.000}{-22.840}{10.9091}{因}\RT{10.909}{-22.840}{10.9091}{此}\RG{132.459}{-25.964}{10.9092}{ex}{32}\RG{172.033}{-25.964}{10.9092}{ex}{33}\RG{141.095}{-37.339}{10.9091}{mi}{117}\RG{161.823}{-37.339}{10.9091}{mi}{118}\RG{291.246}{-40.114}{7.9701}{rm}{50}\RG{319.800}{-40.114}{7.9701}{rm}{50}\RG{343.545}{-40.114}{7.9701}{rm}{50}\RG{109.336}{-44.619}{10.9091}{mi}{71}\RG{120.944}{-44.619}{10.9091}{rm}{61}\RG{182.487}{-44.619}{10.9091}{mi}{59}\RT{209.155}{-44.619}{10.9091}{d}\RT{215.220}{-44.619}{10.9091}{e}\RT{220.064}{-44.619}{10.9091}{t}\RG{224.307}{-44.619}{10.9091}{rm}{40}\RG{228.550}{-44.619}{10.9091}{mi}{120}\RG{234.779}{-44.619}{10.9091}{mi}{73}\RG{246.737}{-44.619}{10.9091}{sy}{0}\RG{257.646}{-44.619}{10.9091}{mi}{71}\RG{266.223}{-44.619}{10.9091}{rm}{41}\RG{273.478}{-44.619}{10.9091}{rm}{61}\RG{285.011}{-44.619}{10.9091}{mi}{120}\RG{298.402}{-44.619}{10.9091}{sy}{0}\RG{309.312}{-44.619}{10.9091}{rm}{40}\RG{313.555}{-44.619}{10.9091}{mi}{117}\RG{326.956}{-44.619}{10.9091}{rm}{43}\RG{337.865}{-44.619}{10.9091}{mi}{118}\RG{348.277}{-44.619}{10.9091}{rm}{41}\RG{352.519}{-44.619}{10.9091}{mi}{58}\RG{239.580}{-46.255}{7.9701}{rm}{50}\RG{141.378}{-52.514}{10.9091}{mi}{118}\RG{157.303}{-52.514}{10.9091}{sy}{0}\RG{165.788}{-52.514}{10.9091}{mi}{117}%
}

由于\allowbreak{}\hskip.1em \RB{28.017}{8.313}{2.854}{%
\RG{0.000}{0.002}{10.9091}{mi}{32}\RG{10.760}{0.002}{10.9091}{sy}{54}\RG{10.760}{0.002}{10.9091}{rm}{61}\RG{22.498}{0.002}{10.9091}{rm}{48}%
}\hskip.1em\allowbreak{} ，\allowbreak{}\hskip.1em \RB{16.678}{5.324}{2.603}{%
\RG{0.000}{0.002}{10.9091}{mi}{117}\RG{6.240}{0.002}{10.9091}{mi}{59}\RG{11.073}{0.002}{10.9091}{mi}{118}%
}\hskip.1em\allowbreak{} 不全为零，故\allowbreak{}\hskip.1em \RB{8.791}{8.193}{0.738}{%
\RG{0.000}{0.002}{10.9091}{mi}{71}%
}\hskip.1em\allowbreak{} 有两个实特征值\allowbreak{}\hskip.1em \RB{51.867}{10.136}{1.773}{%
\RG{0.000}{0.002}{10.9091}{sy}{6}\RG{8.485}{9.200}{10.9091}{sy}{112}\RG{17.576}{0.002}{10.9091}{mi}{117}\RG{23.821}{3.153}{7.9701}{rm}{50}\RG{30.977}{0.002}{10.9091}{rm}{43}\RG{41.886}{0.002}{10.9091}{mi}{118}\RG{47.565}{3.153}{7.9701}{rm}{50}%
}\hskip.1em\allowbreak{} 。取其任一非零实特征向量，所张成的一维子空间就是非平凡的\allowbreak{}\hskip.1em \RB{7.427}{8.068}{2.739}{%
\RG{0.000}{-0.003}{10.9091}{mi}{32}%
}\hskip.1em\allowbreak{} –不变子空间。小结：在二维实空间中，无实特征值与反交换条件共同迫使第一个变换满足\allowbreak{}\hskip.1em \RB{56.617}{9.763}{2.874}{%
\RG{0.000}{0.004}{10.9091}{mi}{39}\RG{7.137}{3.963}{7.9701}{rm}{50}\RG{14.947}{0.004}{10.9091}{rm}{61}\RG{26.511}{0.004}{10.9091}{sy}{0}\RG{34.996}{0.004}{10.9091}{mi}{115}\RG{40.109}{3.963}{7.9701}{rm}{50}\RG{44.842}{0.004}{10.9091}{mi}{73}\RG{49.637}{-1.684}{7.9701}{mi}{86}%
}\hskip.1em\allowbreak{} ，从而第二个变换具有实特征值。

\[
\RB{278.451}{7.724}{2.296}{%
\RT{0.000}{-0.003}{7.9701}{题}\RT{8.966}{-0.003}{7.9701}{目}\RT{17.933}{-0.003}{7.9701}{来}\RT{26.899}{-0.003}{7.9701}{源}\RT{35.866}{-0.003}{7.9701}{：}\RT{44.832}{-0.003}{7.9701}{谢}\RT{53.799}{-0.003}{7.9701}{启}\RT{62.765}{-0.003}{7.9701}{鸿}\RT{71.732}{-0.003}{7.9701}{《}\RT{80.698}{-0.003}{7.9701}{复}\RT{89.665}{-0.003}{7.9701}{旦}\RT{98.631}{-0.003}{7.9701}{大}\RT{107.598}{-0.003}{7.9701}{学}\RT{116.564}{-0.003}{7.9701}{高}\RT{125.531}{-0.003}{7.9701}{等}\RT{134.497}{-0.003}{7.9701}{代}\RT{143.464}{-0.003}{7.9701}{数}\RT{152.430}{-0.003}{7.9701}{习}\RT{161.397}{-0.003}{7.9701}{题}\RT{170.363}{-0.003}{7.9701}{集}\RT{179.329}{-0.003}{7.9701}{（}\RT{188.296}{-0.003}{7.9701}{2}\RT{192.779}{-0.003}{7.9701}{0}\RT{197.262}{-0.003}{7.9701}{2}\RT{201.745}{-0.003}{7.9701}{3}\RT{209.214}{-0.003}{7.9701}{版}\RT{218.180}{-0.003}{7.9701}{）}\RT{222.672}{-0.003}{7.9701}{》}\RT{227.156}{-0.003}{7.9701}{，}\RT{236.122}{-0.003}{7.9701}{第}\RT{248.066}{-0.003}{7.9701}{7}\RT{252.549}{-0.003}{7.9701}{4}\RT{260.018}{-0.003}{7.9701}{页}\RT{268.985}{-0.003}{7.9701}{。}%
}
\]
% Original page 3.
\section*{2019S03}
设\allowbreak{}\hskip.1em \RB{26.491}{7.768}{1.993}{%
\RG{0.000}{0.002}{10.9091}{mi}{110}\RG{9.578}{0.002}{10.9091}{sy}{21}\RG{21.093}{0.002}{10.9091}{rm}{50}%
}\hskip.1em\allowbreak{} ，\allowbreak{}\hskip.1em \RB{36.603}{8.182}{2.603}{%
\RG{0.000}{0.002}{10.9091}{mi}{97}\RG{5.760}{0.002}{10.9091}{mi}{59}\RG{10.593}{0.002}{10.9091}{mi}{98}\RG{18.327}{0.002}{10.9091}{sy}{50}\RG{28.630}{0.002}{10.9091}{bb}{67}%
}\hskip.1em\allowbreak{} ，考察矩阵

\[
\RB{132.161}{11.234}{74.837}{%
\RG{22.727}{10.308}{10.9092}{ex}{48}\RG{118.203}{10.308}{10.9092}{ex}{49}\RG{32.273}{-0.005}{10.9091}{rm}{48}\RG{47.690}{-0.005}{10.9091}{mi}{97}\RG{67.796}{-0.005}{10.9091}{mi}{97}\RG{87.928}{-0.005}{10.9091}{sy}{1}\RG{92.761}{-0.005}{10.9091}{sy}{1}\RG{97.604}{-0.005}{10.9091}{sy}{1}\RG{112.436}{-0.005}{10.9091}{mi}{97}\RG{22.727}{-8.892}{10.9092}{ex}{66}\RG{118.203}{-8.892}{10.9092}{ex}{67}\RG{32.659}{-15.180}{10.9091}{mi}{98}\RG{47.846}{-15.180}{10.9091}{rm}{48}\RG{67.810}{-15.180}{10.9091}{mi}{97}\RG{87.928}{-15.180}{10.9091}{sy}{1}\RG{92.761}{-15.180}{10.9091}{sy}{1}\RG{97.604}{-15.180}{10.9091}{sy}{1}\RG{112.436}{-15.180}{10.9091}{mi}{97}\RG{22.727}{-15.438}{10.9092}{ex}{66}\RG{118.203}{-15.438}{10.9092}{ex}{67}\RG{22.727}{-21.983}{10.9092}{ex}{66}\RG{118.203}{-21.983}{10.9092}{ex}{67}\RT{113.803}{-26.866}{10.9091}{.}\RT{89.440}{-27.863}{10.9091}{.}\RG{22.727}{-28.529}{10.9092}{ex}{66}\RG{118.203}{-28.529}{10.9092}{ex}{67}\RT{93.684}{-30.852}{10.9091}{.}\RT{113.803}{-30.852}{10.9091}{.}\RT{97.929}{-33.840}{10.9091}{.}\RG{0.000}{-34.529}{10.9091}{mi}{65}\RG{11.212}{-34.529}{10.9091}{rm}{61}\RG{129.566}{-34.529}{10.9091}{mi}{58}\RG{32.659}{-34.837}{10.9091}{mi}{98}\RG{48.227}{-34.837}{10.9091}{mi}{98}\RG{67.966}{-34.837}{10.9091}{rm}{48}\RT{113.803}{-34.837}{10.9091}{.}\RG{22.727}{-35.074}{10.9092}{ex}{66}\RG{118.203}{-35.074}{10.9092}{ex}{67}\RG{22.727}{-41.620}{10.9092}{ex}{66}\RG{118.203}{-41.620}{10.9092}{ex}{67}\RT{33.484}{-46.523}{10.9091}{.}\RT{49.057}{-46.523}{10.9091}{.}\RT{64.932}{-47.519}{10.9091}{.}\RT{89.440}{-47.519}{10.9091}{.}\RG{22.727}{-48.165}{10.9092}{ex}{66}\RG{118.203}{-48.165}{10.9092}{ex}{67}\RT{33.484}{-50.508}{10.9091}{.}\RT{49.057}{-50.508}{10.9091}{.}\RT{69.176}{-50.508}{10.9091}{.}\RT{93.684}{-50.508}{10.9091}{.}\RT{73.421}{-53.497}{10.9091}{.}\RT{97.929}{-53.497}{10.9091}{.}\RT{33.484}{-54.493}{10.9091}{.}\RT{49.057}{-54.493}{10.9091}{.}\RG{112.436}{-54.493}{10.9091}{mi}{97}\RG{22.727}{-55.147}{10.9092}{ex}{64}\RG{118.203}{-55.147}{10.9092}{ex}{65}\RG{32.659}{-69.669}{10.9091}{mi}{98}\RG{48.227}{-69.669}{10.9091}{mi}{98}\RG{63.419}{-69.669}{10.9091}{sy}{1}\RG{68.252}{-69.669}{10.9091}{sy}{1}\RG{73.096}{-69.669}{10.9091}{sy}{1}\RG{92.860}{-69.669}{10.9091}{mi}{98}\RG{112.592}{-69.669}{10.9091}{rm}{48}%
}
\]
求\allowbreak{}\hskip.1em \RB{8.366}{8.315}{0.500}{%
\RG{0.000}{0.004}{10.9091}{mi}{65}%
}\hskip.1em\allowbreak{} 可对角化的充分必要条件。

解：所求条件为

\allowbreak{}\hskip.1em \RB{45.041}{8.068}{0.743}{%
\RG{0.000}{-0.003}{10.9091}{mi}{97}\RG{8.797}{-0.003}{10.9091}{rm}{61}\RG{20.312}{-0.003}{10.9091}{mi}{98}\RG{28.024}{-0.003}{10.9091}{rm}{61}\RG{39.523}{-0.003}{10.9091}{rm}{48}%
}\hskip.1em\allowbreak{} 或\allowbreak{}\hskip.1em \RB{33.043}{8.308}{2.859}{%
\RG{0.000}{-0.003}{10.9091}{mi}{97}\RG{5.760}{-0.003}{10.9091}{mi}{98}\RG{13.478}{-0.003}{10.9091}{sy}{54}\RG{13.478}{-0.003}{10.9091}{rm}{61}\RG{24.976}{-0.003}{10.9091}{rm}{48}\RG{30.448}{-0.003}{10.9091}{mi}{58}%
}\hskip.1em\allowbreak{}

先证充分性。若\allowbreak{}\hskip.1em \RB{45.041}{8.070}{0.741}{%
\RG{0.000}{-0.001}{10.9091}{mi}{97}\RG{8.797}{-0.001}{10.9091}{rm}{61}\RG{20.312}{-0.001}{10.9091}{mi}{98}\RG{28.024}{-0.001}{10.9091}{rm}{61}\RG{39.523}{-0.001}{10.9091}{rm}{48}%
}\hskip.1em\allowbreak{} ，则\allowbreak{}\hskip.1em \RB{28.229}{8.310}{0.741}{%
\RG{0.000}{-0.001}{10.9091}{mi}{65}\RG{11.212}{-0.001}{10.9091}{rm}{61}\RG{22.710}{-0.001}{10.9091}{rm}{48}%
}\hskip.1em\allowbreak{} 。若\allowbreak{}\hskip.1em \RB{45.040}{8.310}{2.857}{%
\RG{0.000}{-0.001}{10.9091}{mi}{97}\RG{8.796}{-0.001}{10.9091}{rm}{61}\RG{20.311}{-0.001}{10.9091}{mi}{98}\RG{28.023}{-0.001}{10.9091}{sy}{54}\RG{28.023}{-0.001}{10.9091}{rm}{61}\RG{39.522}{-0.001}{10.9091}{rm}{48}%
}\hskip.1em\allowbreak{} ，记\allowbreak{}\hskip.1em \RB{7.406}{7.950}{0.741}{%
\RG{0.000}{-0.001}{10.9091}{mi}{74}%
}\hskip.1em\allowbreak{} 为所有元素均为\allowbreak{}\hskip.1em \RB{5.071}{7.765}{0.501}{%
\RG{0.000}{-0.001}{10.9091}{rm}{49}%
}\hskip.1em\allowbreak{} 的\allowbreak{}\hskip.1em \RB{6.729}{5.321}{0.621}{%
\RG{0.000}{-0.001}{10.9091}{mi}{110}%
}\hskip.1em\allowbreak{} 阶矩阵，则

\[
\RB{70.437}{8.677}{3.232}{%
\RG{0.000}{-0.005}{10.9091}{mi}{65}\RG{11.212}{-0.005}{10.9091}{rm}{61}\RG{22.727}{-0.005}{10.9091}{mi}{97}\RG{28.493}{-0.005}{10.9091}{rm}{40}\RG{32.736}{-0.005}{10.9091}{mi}{74}\RG{42.258}{-0.005}{10.9091}{sy}{0}\RG{53.167}{-0.005}{10.9091}{mi}{73}\RG{63.600}{-0.005}{10.9091}{rm}{41}\RG{67.843}{-0.005}{10.9091}{mi}{58}\RG{57.963}{-1.641}{7.9701}{mi}{110}%
}
\]
由于\allowbreak{}\hskip.1em \RB{7.406}{7.952}{0.739}{%
\RG{0.000}{0.001}{10.9091}{mi}{74}%
}\hskip.1em\allowbreak{} 在span\allowbreak{}\hskip.1em \RB{48.502}{8.683}{3.226}{%
\RG{0.000}{0.001}{10.9091}{sy}{102}\RG{5.455}{0.001}{10.9091}{rm}{40}\RG{9.688}{0.001}{10.9091}{rm}{49}\RG{15.152}{0.001}{10.9091}{mi}{59}\RG{19.985}{0.001}{10.9091}{mi}{58}\RG{24.829}{0.001}{10.9091}{mi}{58}\RG{29.661}{0.001}{10.9091}{mi}{58}\RG{34.505}{0.001}{10.9091}{mi}{59}\RG{39.394}{0.001}{10.9091}{rm}{49}\RG{44.849}{0.001}{10.9091}{rm}{41}%
}\hskip.1em\allowbreak{} T\allowbreak{}\hskip.1em \RB{5.158}{8.683}{3.226}{%
\RG{0.000}{0.001}{10.9091}{sy}{103}%
}\hskip.1em\allowbreak{} 上的特征值为\allowbreak{}\hskip.1em \RB{6.729}{5.323}{0.619}{%
\RG{0.000}{0.001}{10.9091}{mi}{110}%
}\hskip.1em\allowbreak{} ，在坐标和为零的\allowbreak{}\hskip.1em \RB{25.550}{7.767}{0.619}{%
\RG{0.000}{0.001}{10.9091}{mi}{110}\RG{9.271}{0.001}{10.9091}{sy}{0}\RG{20.479}{0.001}{10.9091}{rm}{49}%
}\hskip.1em\allowbreak{} 维子空间上的特征值为\allowbreak{}\hskip.1em \RB{5.518}{7.767}{0.739}{%
\RG{0.000}{0.001}{10.9091}{rm}{48}%
}\hskip.1em\allowbreak{} ，故\allowbreak{}\hskip.1em \RB{8.366}{8.307}{0.504}{%
\RG{0.000}{-0.004}{10.9091}{mi}{65}%
}\hskip.1em\allowbreak{} 可对角化。

若\allowbreak{}\hskip.1em \RB{30.495}{8.315}{2.852}{%
\RG{0.000}{0.004}{10.9091}{mi}{97}\RG{5.760}{0.004}{10.9091}{mi}{98}\RG{13.478}{0.004}{10.9091}{sy}{54}\RG{13.478}{0.004}{10.9091}{rm}{61}\RG{24.976}{0.004}{10.9091}{rm}{48}%
}\hskip.1em\allowbreak{} 且\allowbreak{}\hskip.1em \RB{25.338}{8.315}{2.852}{%
\RG{0.000}{0.004}{10.9091}{mi}{97}\RG{8.796}{0.004}{10.9091}{sy}{54}\RG{8.796}{0.004}{10.9091}{rm}{61}\RG{20.311}{0.004}{10.9091}{mi}{98}%
}\hskip.1em\allowbreak{} ，令\allowbreak{}\hskip.1em \RB{4.100}{7.333}{0.616}{%
\RG{0.000}{0.004}{10.9091}{mi}{116}%
}\hskip.1em\allowbreak{} 遍历方程

\[
\RB{31.249}{15.452}{8.103}{%
\RG{25.858}{7.381}{10.9091}{mi}{98}\RG{3.939}{4.505}{7.9701}{mi}{110}\RG{0.000}{-0.000}{10.9091}{mi}{116}\RG{12.606}{-0.000}{10.9091}{rm}{61}\RG{25.316}{-7.483}{10.9091}{mi}{97}\RR{25.316}{2.728}{5.767}{0.436}%
}
\]
的\allowbreak{}\hskip.1em \RB{6.729}{5.322}{0.620}{%
\RG{0.000}{0.000}{10.9091}{mi}{110}%
}\hskip.1em\allowbreak{} 个不同复根。每个根都满足\allowbreak{}\hskip.1em \RB{33.859}{8.311}{2.856}{%
\RG{0.000}{0.000}{10.9091}{mi}{116}\RG{6.970}{0.000}{10.9091}{sy}{54}\RG{6.970}{0.000}{10.9091}{rm}{61}\RG{18.468}{0.000}{10.9091}{rm}{48}\RG{23.940}{0.000}{10.9091}{mi}{59}\RG{28.788}{0.000}{10.9091}{rm}{49}%
}\hskip.1em\allowbreak{} 。置

\[
\RB{199.007}{15.449}{8.105}{%
\RG{167.496}{7.378}{10.9091}{mi}{97}\RG{173.256}{7.378}{10.9091}{mi}{116}\RG{179.626}{7.378}{10.9091}{sy}{0}\RG{190.536}{7.378}{10.9091}{mi}{98}\RG{79.317}{4.502}{7.9701}{mi}{110}\RG{84.455}{4.502}{7.9701}{sy}{0}\RG{91.041}{4.502}{7.9701}{rm}{49}\RT{100.016}{4.502}{7.9701}{T}\RG{0.000}{-0.002}{10.9091}{mi}{118}\RG{5.679}{-0.002}{10.9091}{rm}{40}\RG{9.922}{-0.002}{10.9091}{mi}{116}\RG{13.861}{-0.002}{10.9091}{rm}{41}\RG{21.116}{-0.002}{10.9091}{rm}{61}\RG{32.625}{-0.002}{10.9091}{rm}{40}\RG{36.858}{-0.002}{10.9091}{rm}{49}\RG{42.346}{-0.002}{10.9091}{mi}{59}\RG{47.179}{-0.002}{10.9091}{mi}{116}\RG{51.117}{-0.002}{10.9091}{mi}{59}\RG{55.961}{-0.002}{10.9091}{mi}{58}\RG{60.794}{-0.002}{10.9091}{mi}{58}\RG{65.637}{-0.002}{10.9091}{mi}{58}\RG{70.481}{-0.002}{10.9091}{mi}{59}\RG{75.314}{-0.002}{10.9091}{mi}{116}\RG{95.774}{-0.002}{10.9091}{rm}{41}\RG{106.300}{-0.002}{10.9091}{mi}{59}\RG{132.951}{-0.002}{10.9091}{mi}{21}\RG{139.331}{-0.002}{10.9091}{rm}{40}\RG{143.573}{-0.002}{10.9091}{mi}{116}\RG{147.513}{-0.002}{10.9091}{rm}{41}\RG{154.768}{-0.002}{10.9091}{rm}{61}\RG{196.413}{-0.002}{10.9091}{mi}{58}\RG{169.993}{-7.485}{10.9091}{rm}{49}\RG{177.871}{-7.485}{10.9091}{sy}{0}\RG{188.782}{-7.485}{10.9091}{mi}{116}\RR{167.496}{2.725}{27.721}{0.436}%
}
\]
对任意\allowbreak{}\hskip.1em \RB{45.033}{7.770}{1.991}{%
\RG{0.000}{0.004}{10.9091}{rm}{49}\RG{8.485}{0.004}{10.9091}{sy}{20}\RG{20.000}{0.004}{10.9091}{mi}{105}\RG{26.788}{0.004}{10.9091}{sy}{20}\RG{38.304}{0.004}{10.9091}{mi}{110}%
}\hskip.1em\allowbreak{} ，利用等比数列求和公式及\allowbreak{}\hskip.1em \RB{34.916}{8.075}{0.616}{%
\RG{0.000}{0.004}{10.9091}{mi}{97}\RG{5.760}{0.004}{10.9091}{mi}{116}\RG{9.706}{3.963}{7.9701}{mi}{110}\RG{18.373}{0.004}{10.9091}{rm}{61}\RG{29.888}{0.004}{10.9091}{mi}{98}%
}\hskip.1em\allowbreak{} ，有

\[
\RB{161.299}{19.434}{71.761}{%
\RG{60.223}{13.634}{7.9701}{mi}{105}\RG{63.108}{13.634}{7.9701}{sy}{0}\RG{69.694}{13.634}{7.9701}{rm}{49}\RG{126.352}{13.634}{7.9701}{mi}{110}\RG{59.197}{10.361}{10.9092}{ex}{88}\RG{121.042}{10.361}{10.9092}{ex}{88}\RG{80.713}{4.501}{7.9701}{mi}{106}\RG{84.597}{4.501}{7.9701}{sy}{0}\RG{91.183}{4.501}{7.9701}{rm}{49}\RG{146.765}{4.501}{7.9701}{mi}{106}\RG{150.649}{4.501}{7.9701}{sy}{0}\RG{157.236}{4.501}{7.9701}{rm}{49}\RG{0.000}{-0.003}{10.9091}{rm}{40}\RG{4.242}{-0.003}{10.9091}{mi}{65}\RG{12.424}{-0.003}{10.9091}{mi}{118}\RG{18.103}{-0.003}{10.9091}{rm}{40}\RG{22.346}{-0.003}{10.9091}{mi}{116}\RG{26.285}{-0.003}{10.9091}{rm}{41}\RG{30.518}{-0.003}{10.9091}{rm}{41}\RG{41.182}{-0.003}{10.9091}{rm}{61}\RG{52.697}{-0.003}{10.9091}{mi}{98}\RG{76.773}{-0.003}{10.9091}{mi}{116}\RG{98.340}{-0.003}{10.9091}{rm}{43}\RG{109.249}{-0.003}{10.9091}{mi}{97}\RG{142.826}{-0.003}{10.9091}{mi}{116}\RG{34.770}{-1.639}{7.9701}{mi}{105}\RG{59.723}{-13.090}{7.9701}{mi}{106}\RG{63.608}{-13.090}{7.9701}{rm}{61}\RG{70.191}{-13.090}{7.9701}{rm}{49}\RG{116.833}{-13.090}{7.9701}{mi}{106}\RG{120.717}{-13.090}{7.9701}{rm}{61}\RG{127.304}{-13.090}{7.9701}{mi}{105}\RG{130.187}{-13.090}{7.9701}{rm}{43}\RG{136.770}{-13.090}{7.9701}{rm}{49}\RG{85.545}{-24.985}{7.9701}{mi}{105}\RG{88.428}{-24.985}{7.9701}{sy}{0}\RG{95.014}{-24.985}{7.9701}{rm}{49}\RG{131.271}{-24.985}{7.9701}{mi}{105}\RG{151.925}{-24.985}{7.9701}{mi}{110}\RG{53.892}{-28.944}{10.9091}{mi}{98}\RG{58.574}{-28.944}{10.9091}{rm}{40}\RG{62.807}{-28.944}{10.9091}{rm}{49}\RG{70.696}{-28.944}{10.9091}{sy}{0}\RG{81.605}{-28.944}{10.9091}{mi}{116}\RG{99.747}{-28.944}{10.9091}{rm}{41}\RG{106.401}{-28.944}{10.9091}{rm}{43}\RG{117.322}{-28.944}{10.9091}{mi}{97}\RG{123.089}{-28.944}{10.9091}{rm}{40}\RG{127.331}{-28.944}{10.9091}{mi}{116}\RG{137.076}{-28.944}{10.9091}{sy}{0}\RG{147.985}{-28.944}{10.9091}{mi}{116}\RG{157.561}{-28.944}{10.9091}{rm}{41}\RG{41.182}{-36.324}{10.9091}{rm}{61}\RG{96.485}{-43.807}{10.9091}{rm}{49}\RG{104.363}{-43.807}{10.9091}{sy}{0}\RG{115.272}{-43.807}{10.9091}{mi}{116}\RG{53.892}{-56.277}{10.9091}{mi}{97}\RG{59.653}{-56.277}{10.9091}{mi}{116}\RG{66.023}{-56.277}{10.9091}{sy}{0}\RG{76.932}{-56.277}{10.9091}{mi}{98}\RG{88.567}{-59.152}{7.9701}{mi}{105}\RG{91.450}{-59.152}{7.9701}{sy}{0}\RG{98.037}{-59.152}{7.9701}{rm}{49}\RG{41.182}{-63.658}{10.9091}{rm}{61}\RG{84.628}{-63.658}{10.9091}{mi}{116}\RG{105.799}{-63.658}{10.9091}{rm}{61}\RG{117.315}{-63.658}{10.9091}{mi}{21}\RG{123.678}{-63.658}{10.9091}{rm}{40}\RG{127.921}{-63.658}{10.9091}{mi}{116}\RG{131.860}{-63.658}{10.9091}{rm}{41}\RG{136.102}{-63.658}{10.9091}{mi}{118}\RG{141.782}{-63.658}{10.9091}{rm}{40}\RG{146.024}{-63.658}{10.9091}{mi}{116}\RG{149.964}{-63.658}{10.9091}{rm}{41}\RG{157.587}{-63.658}{10.9091}{mi}{58}\RG{154.206}{-65.294}{7.9701}{mi}{105}\RG{56.390}{-71.141}{10.9091}{rm}{49}\RG{64.269}{-71.141}{10.9091}{sy}{0}\RG{75.178}{-71.141}{10.9091}{mi}{116}\RR{53.892}{-33.596}{107.911}{0.436}\RR{53.892}{-60.930}{27.721}{0.436}%
}
\]
故\allowbreak{}\hskip.1em \RB{17.514}{8.678}{3.231}{%
\RG{0.000}{-0.004}{10.9091}{mi}{118}\RG{5.679}{-0.004}{10.9091}{rm}{40}\RG{9.922}{-0.004}{10.9091}{mi}{116}\RG{13.861}{-0.004}{10.9091}{rm}{41}%
}\hskip.1em\allowbreak{} 是\allowbreak{}\hskip.1em \RB{8.366}{8.307}{0.504}{%
\RG{0.000}{-0.004}{10.9091}{mi}{65}%
}\hskip.1em\allowbreak{} 的特征向量。不同根所对应的\allowbreak{}\hskip.1em \RB{6.729}{5.318}{0.624}{%
\RG{0.000}{-0.004}{10.9091}{mi}{110}%
}\hskip.1em\allowbreak{} 个向量构成Vandermonde矩阵的列，其行列式是非零的根差乘积，因此线性无关，\allowbreak{}\hskip.1em \RB{8.366}{8.312}{0.500}{%
\RG{0.000}{0.001}{10.9091}{mi}{65}%
}\hskip.1em\allowbreak{} 可对角化。

再证必要性。若\allowbreak{}\hskip.1em \RB{30.496}{8.070}{0.741}{%
\RG{0.000}{-0.001}{10.9091}{mi}{97}\RG{5.760}{-0.001}{10.9091}{mi}{98}\RG{13.479}{-0.001}{10.9091}{rm}{61}\RG{24.977}{-0.001}{10.9091}{rm}{48}%
}\hskip.1em\allowbreak{} 而\allowbreak{}\hskip.1em \RB{61.981}{8.681}{3.228}{%
\RG{0.000}{-0.001}{10.9091}{rm}{40}\RG{4.243}{-0.001}{10.9091}{mi}{97}\RG{10.003}{-0.001}{10.9091}{mi}{59}\RG{14.836}{-0.001}{10.9091}{mi}{98}\RG{19.539}{-0.001}{10.9091}{rm}{41}\RG{26.812}{-0.001}{10.9091}{sy}{54}\RG{26.812}{-0.001}{10.9091}{rm}{61}\RG{38.310}{-0.001}{10.9091}{rm}{40}\RG{42.543}{-0.001}{10.9091}{rm}{48}\RG{48.024}{-0.001}{10.9091}{mi}{59}\RG{52.873}{-0.001}{10.9091}{rm}{48}\RG{58.328}{-0.001}{10.9091}{rm}{41}%
}\hskip.1em\allowbreak{} ，则\allowbreak{}\hskip.1em \RB{8.366}{8.310}{0.501}{%
\RG{0.000}{-0.001}{10.9091}{mi}{65}%
}\hskip.1em\allowbreak{} 是非零严格上三角阵或严格下三角阵，因而是非零幂零阵。若其可对角化，则由于所有特征值均为\allowbreak{}\hskip.1em \RB{5.518}{7.770}{0.736}{%
\RG{0.000}{0.004}{10.9091}{rm}{48}%
}\hskip.1em\allowbreak{} ，必有\allowbreak{}\hskip.1em \RB{28.229}{8.315}{0.736}{%
\RG{0.000}{0.004}{10.9091}{mi}{65}\RG{11.212}{0.004}{10.9091}{rm}{61}\RG{22.710}{0.004}{10.9091}{rm}{48}%
}\hskip.1em\allowbreak{} ，矛盾。

综上，得到所求充要条件。小结：当\allowbreak{}\hskip.1em \RB{30.709}{8.310}{2.857}{%
\RG{0.000}{-0.001}{10.9091}{mi}{97}\RG{5.760}{-0.001}{10.9091}{mi}{98}\RG{13.583}{-0.001}{10.9091}{sy}{54}\RG{13.583}{-0.001}{10.9091}{rm}{61}\RG{25.190}{-0.001}{10.9091}{rm}{48}%
}\hskip.1em\allowbreak{} 时，通过\allowbreak{}\hskip.1em \RB{40.398}{8.681}{3.228}{%
\RG{0.000}{-0.001}{10.9091}{mi}{116}\RG{3.940}{3.958}{7.9701}{mi}{110}\RG{12.710}{-0.001}{10.9091}{rm}{61}\RG{24.329}{-0.001}{10.9091}{mi}{98}\RG{29.011}{-0.001}{10.9091}{rm}{47}\RG{34.465}{-0.001}{10.9091}{mi}{97}%
}\hskip.1em\allowbreak{} 构造Vandermonde型特征向量；恰有一个参数为零时出现非零幂零阵。

\[
\RB{278.451}{7.729}{2.291}{%
\RT{0.000}{0.002}{7.9701}{题}\RT{8.966}{0.002}{7.9701}{目}\RT{17.933}{0.002}{7.9701}{来}\RT{26.899}{0.002}{7.9701}{源}\RT{35.866}{0.002}{7.9701}{：}\RT{44.832}{0.002}{7.9701}{谢}\RT{53.799}{0.002}{7.9701}{启}\RT{62.765}{0.002}{7.9701}{鸿}\RT{71.732}{0.002}{7.9701}{《}\RT{80.698}{0.002}{7.9701}{复}\RT{89.665}{0.002}{7.9701}{旦}\RT{98.631}{0.002}{7.9701}{大}\RT{107.598}{0.002}{7.9701}{学}\RT{116.564}{0.002}{7.9701}{高}\RT{125.531}{0.002}{7.9701}{等}\RT{134.497}{0.002}{7.9701}{代}\RT{143.464}{0.002}{7.9701}{数}\RT{152.430}{0.002}{7.9701}{习}\RT{161.397}{0.002}{7.9701}{题}\RT{170.363}{0.002}{7.9701}{集}\RT{179.329}{0.002}{7.9701}{（}\RT{188.296}{0.002}{7.9701}{2}\RT{192.779}{0.002}{7.9701}{0}\RT{197.262}{0.002}{7.9701}{2}\RT{201.745}{0.002}{7.9701}{3}\RT{209.214}{0.002}{7.9701}{版}\RT{218.180}{0.002}{7.9701}{）}\RT{222.672}{0.002}{7.9701}{》}\RT{227.156}{0.002}{7.9701}{，}\RT{236.122}{0.002}{7.9701}{第}\RT{248.066}{0.002}{7.9701}{7}\RT{252.549}{0.002}{7.9701}{4}\RT{260.018}{0.002}{7.9701}{页}\RT{268.985}{0.002}{7.9701}{。}%
}
\]
% Original page 4.
\section*{2019S04}
设\allowbreak{}\hskip.1em \RB{21.762}{8.313}{2.603}{%
\RG{0.000}{0.002}{10.9091}{mi}{65}\RG{8.182}{0.002}{10.9091}{mi}{59}\RG{13.015}{0.002}{10.9091}{mi}{66}%
}\hskip.1em\allowbreak{} 为\allowbreak{}\hskip.1em \RB{6.729}{5.324}{0.618}{%
\RG{0.000}{0.002}{10.9091}{mi}{110}%
}\hskip.1em\allowbreak{} 阶方阵，满足

\[
\RB{98.188}{10.301}{1.408}{%
\RG{8.182}{4.501}{7.9701}{rm}{50}\RG{70.862}{4.501}{7.9701}{rm}{50}\RG{0.000}{-0.003}{10.9091}{mi}{65}\RG{15.338}{-0.003}{10.9091}{sy}{0}\RG{26.248}{-0.003}{10.9091}{rm}{50}\RG{31.702}{-0.003}{10.9091}{mi}{65}\RG{39.884}{-0.003}{10.9091}{mi}{66}\RG{51.130}{-0.003}{10.9091}{rm}{43}\RG{62.039}{-0.003}{10.9091}{mi}{66}\RG{78.625}{-0.003}{10.9091}{rm}{61}\RG{90.123}{-0.003}{10.9091}{rm}{48}\RG{95.594}{-0.003}{10.9091}{mi}{58}%
}
\]
（1）若\allowbreak{}\hskip.1em \RB{26.475}{7.762}{0.624}{%
\RG{0.000}{-0.004}{10.9091}{mi}{110}\RG{9.578}{-0.004}{10.9091}{rm}{61}\RG{21.076}{-0.004}{10.9091}{rm}{50}%
}\hskip.1em\allowbreak{} ，证明：\allowbreak{}\hskip.1em \RB{48.730}{8.307}{0.504}{%
\RG{0.000}{-0.004}{10.9091}{mi}{65}\RG{8.182}{-0.004}{10.9091}{mi}{66}\RG{20.035}{-0.004}{10.9091}{rm}{61}\RG{31.550}{-0.004}{10.9091}{mi}{66}\RG{40.365}{-0.004}{10.9091}{mi}{65}%
}\hskip.1em\allowbreak{} 。

（2）若\allowbreak{}\hskip.1em \RB{26.580}{7.770}{1.991}{%
\RG{0.000}{0.004}{10.9091}{mi}{110}\RG{9.579}{0.004}{10.9091}{sy}{21}\RG{21.094}{0.004}{10.9091}{rm}{51}%
}\hskip.1em\allowbreak{} ，举例说明：\allowbreak{}\hskip.1em \RB{48.729}{8.315}{0.500}{%
\RG{0.000}{0.004}{10.9091}{mi}{65}\RG{8.182}{0.004}{10.9091}{mi}{66}\RG{20.034}{0.004}{10.9091}{rm}{61}\RG{31.549}{0.004}{10.9091}{mi}{66}\RG{40.364}{0.004}{10.9091}{mi}{65}%
}\hskip.1em\allowbreak{} 不一定成立。

证明：（1）令\allowbreak{}\hskip.1em \RB{53.386}{8.312}{0.739}{%
\RG{0.000}{0.001}{10.9091}{mi}{67}\RG{11.608}{0.001}{10.9091}{rm}{61}\RG{23.123}{0.001}{10.9091}{mi}{65}\RG{33.729}{0.001}{10.9091}{sy}{0}\RG{44.638}{0.001}{10.9091}{mi}{66}%
}\hskip.1em\allowbreak{} 。整理题设等式得

\[
\RB{139.848}{10.302}{0.743}{%
\RG{8.578}{4.502}{7.9701}{rm}{50}\RG{0.000}{-0.003}{10.9091}{mi}{67}\RG{16.340}{-0.003}{10.9091}{rm}{61}\RG{27.855}{-0.003}{10.9091}{mi}{65}\RG{36.037}{-0.003}{10.9091}{mi}{66}\RG{47.283}{-0.003}{10.9091}{sy}{0}\RG{58.192}{-0.003}{10.9091}{mi}{66}\RG{67.007}{-0.003}{10.9091}{mi}{65}\RG{78.227}{-0.003}{10.9091}{rm}{61}\RG{89.742}{-0.003}{10.9091}{mi}{67}\RG{98.305}{-0.003}{10.9091}{mi}{66}\RG{109.566}{-0.003}{10.9091}{sy}{0}\RG{120.475}{-0.003}{10.9091}{mi}{66}\RG{129.290}{-0.003}{10.9091}{mi}{67}\RG{137.253}{-0.003}{10.9091}{mi}{58}%
}
\]
因为\allowbreak{}\hskip.1em \RB{8.791}{8.194}{0.737}{%
\RG{0.000}{0.003}{10.9091}{mi}{67}%
}\hskip.1em\allowbreak{} 与\allowbreak{}\hskip.1em \RB{8.791}{8.194}{0.737}{%
\RG{0.000}{0.003}{10.9091}{mi}{67}%
}\hskip.1em\allowbreak{} 的任意幂可交换，所以对所有正整数\allowbreak{}\hskip.1em \RB{5.256}{5.325}{0.617}{%
\RG{0.000}{0.003}{10.9091}{mi}{114}%
}\hskip.1em\allowbreak{} ，

\[
\RB{171.995}{19.436}{54.790}{%
\RG{47.849}{13.636}{7.9701}{mi}{114}\RG{51.905}{13.636}{7.9701}{sy}{0}\RG{58.493}{13.636}{7.9701}{rm}{49}\RG{47.408}{10.363}{10.9092}{ex}{88}\RG{11.608}{4.504}{7.9701}{mi}{114}\RG{73.562}{4.504}{7.9701}{mi}{106}\RG{114.224}{4.504}{7.9701}{mi}{114}\RG{118.281}{4.504}{7.9701}{sy}{0}\RG{124.867}{4.504}{7.9701}{rm}{49}\RG{129.101}{4.504}{7.9701}{sy}{0}\RG{135.688}{4.504}{7.9701}{mi}{106}\RG{0.000}{-0.001}{10.9091}{rm}{91}\RG{3.030}{-0.001}{10.9091}{mi}{67}\RG{16.162}{-0.001}{10.9091}{mi}{59}\RG{20.995}{-0.001}{10.9091}{mi}{66}\RG{29.833}{-0.001}{10.9091}{rm}{93}\RG{35.877}{-0.001}{10.9091}{rm}{61}\RG{64.985}{-0.001}{10.9091}{mi}{67}\RG{77.945}{-0.001}{10.9091}{rm}{91}\RG{80.975}{-0.001}{10.9091}{mi}{67}\RG{88.927}{-0.001}{10.9091}{mi}{59}\RG{93.771}{-0.001}{10.9091}{mi}{66}\RG{102.617}{-0.001}{10.9091}{rm}{93}\RG{105.647}{-0.001}{10.9091}{mi}{67}\RG{47.935}{-13.088}{7.9701}{mi}{106}\RG{51.819}{-13.088}{7.9701}{rm}{61}\RG{58.403}{-13.088}{7.9701}{rm}{48}\RG{47.849}{-25.941}{7.9701}{mi}{114}\RG{51.905}{-25.941}{7.9701}{sy}{0}\RG{58.492}{-25.941}{7.9701}{rm}{49}\RG{47.408}{-29.214}{10.9092}{ex}{88}\RG{73.562}{-35.073}{7.9701}{mi}{106}\RG{86.522}{-35.073}{7.9701}{rm}{50}\RG{99.831}{-35.073}{7.9701}{mi}{114}\RG{103.888}{-35.073}{7.9701}{sy}{0}\RG{110.475}{-35.073}{7.9701}{rm}{49}\RG{114.709}{-35.073}{7.9701}{sy}{0}\RG{121.295}{-35.073}{7.9701}{mi}{106}\RG{154.025}{-35.073}{7.9701}{mi}{114}\RG{158.081}{-35.073}{7.9701}{rm}{43}\RG{164.665}{-35.073}{7.9701}{rm}{49}\RG{35.893}{-39.577}{10.9091}{rm}{61}\RG{64.985}{-39.577}{10.9091}{mi}{67}\RG{77.945}{-39.577}{10.9091}{mi}{67}\RG{91.254}{-39.577}{10.9091}{mi}{67}\RG{128.708}{-39.577}{10.9091}{rm}{61}\RG{140.223}{-39.577}{10.9091}{mi}{114}\RG{145.438}{-39.577}{10.9091}{mi}{67}\RG{169.400}{-39.577}{10.9091}{mi}{58}\RG{47.935}{-52.664}{7.9701}{mi}{106}\RG{51.819}{-52.664}{7.9701}{rm}{61}\RG{58.403}{-52.664}{7.9701}{rm}{48}%
}
\]
取迹，并利用交换子的迹为零，得

\[
\RB{97.769}{10.307}{3.224}{%
\RG{21.339}{4.507}{7.9701}{rm}{50}\RG{66.200}{4.507}{7.9701}{rm}{51}\RT{0.000}{0.003}{10.9091}{t}\RT{4.244}{0.003}{10.9091}{r}\RG{8.520}{0.003}{10.9091}{rm}{40}\RG{12.762}{0.003}{10.9091}{mi}{67}\RG{26.072}{0.003}{10.9091}{rm}{41}\RG{33.327}{0.003}{10.9091}{rm}{61}\RT{44.859}{0.003}{10.9091}{t}\RT{49.103}{0.003}{10.9091}{r}\RG{53.379}{0.003}{10.9091}{rm}{40}\RG{57.622}{0.003}{10.9091}{mi}{67}\RG{70.932}{0.003}{10.9091}{rm}{41}\RG{78.186}{0.003}{10.9091}{rm}{61}\RG{89.695}{0.003}{10.9091}{rm}{48}\RG{95.175}{0.003}{10.9091}{mi}{58}%
}
\]
设\allowbreak{}\hskip.1em \RB{15.335}{7.328}{0.621}{%
\RG{0.000}{-0.001}{10.9091}{mi}{116}\RG{6.969}{-0.001}{10.9091}{rm}{61}%
}\hskip.1em\allowbreak{} tr\allowbreak{}\hskip.1em \RB{8.791}{8.190}{0.741}{%
\RG{0.000}{-0.001}{10.9091}{mi}{67}%
}\hskip.1em\allowbreak{} 、\allowbreak{}\hskip.1em \RB{17.074}{8.070}{0.621}{%
\RG{0.000}{-0.001}{10.9091}{mi}{100}\RG{8.708}{-0.001}{10.9091}{rm}{61}%
}\hskip.1em\allowbreak{} det\allowbreak{}\hskip.1em \RB{8.791}{8.190}{0.741}{%
\RG{0.000}{-0.001}{10.9091}{mi}{67}%
}\hskip.1em\allowbreak{} 。由二阶矩阵的特征多项式与迹的关系，

\[
\RB{229.197}{10.299}{3.232}{%
\RG{21.340}{4.499}{7.9701}{rm}{50}\RG{48.800}{4.499}{7.9701}{rm}{50}\RG{146.005}{4.499}{7.9701}{rm}{51}\RG{173.465}{4.499}{7.9701}{rm}{51}\RT{0.000}{-0.005}{10.9091}{t}\RT{4.244}{-0.005}{10.9091}{r}\RG{8.520}{-0.005}{10.9091}{rm}{40}\RG{12.763}{-0.005}{10.9091}{mi}{67}\RG{26.072}{-0.005}{10.9091}{rm}{41}\RG{33.327}{-0.005}{10.9091}{rm}{61}\RG{44.860}{-0.005}{10.9091}{mi}{116}\RG{55.956}{-0.005}{10.9091}{sy}{0}\RG{66.866}{-0.005}{10.9091}{rm}{50}\RG{72.321}{-0.005}{10.9091}{mi}{100}\RG{81.029}{-0.005}{10.9091}{rm}{61}\RG{92.527}{-0.005}{10.9091}{rm}{48}\RG{97.999}{-0.005}{10.9091}{mi}{59}\RT{124.665}{-0.005}{10.9091}{t}\RT{128.909}{-0.005}{10.9091}{r}\RG{133.185}{-0.005}{10.9091}{rm}{40}\RG{137.428}{-0.005}{10.9091}{mi}{67}\RG{150.737}{-0.005}{10.9091}{rm}{41}\RG{157.992}{-0.005}{10.9091}{rm}{61}\RG{169.525}{-0.005}{10.9091}{mi}{116}\RG{180.621}{-0.005}{10.9091}{sy}{0}\RG{191.531}{-0.005}{10.9091}{rm}{51}\RG{196.986}{-0.005}{10.9091}{mi}{116}\RG{200.924}{-0.005}{10.9091}{mi}{100}\RG{209.633}{-0.005}{10.9091}{rm}{61}\RG{221.131}{-0.005}{10.9091}{rm}{48}\RG{226.603}{-0.005}{10.9091}{mi}{58}%
}
\]
于是\allowbreak{}\hskip.1em \RB{45.711}{9.760}{3.226}{%
\RG{0.000}{0.001}{10.9091}{mi}{100}\RG{11.690}{0.001}{10.9091}{rm}{61}\RG{26.186}{0.001}{10.9091}{mi}{116}\RG{30.126}{3.960}{7.9701}{rm}{50}\RG{34.858}{0.001}{10.9091}{rm}{47}\RG{40.313}{0.001}{10.9091}{rm}{50}%
}\hskip.1em\allowbreak{} ，继而\allowbreak{}\hskip.1em \RB{123.908}{9.760}{3.226}{%
\RG{0.000}{0.001}{10.9091}{rm}{48}\RG{11.466}{0.001}{10.9091}{rm}{61}\RG{25.963}{0.001}{10.9091}{mi}{116}\RG{29.902}{3.960}{7.9701}{rm}{51}\RG{38.251}{0.001}{10.9091}{sy}{0}\RG{50.353}{0.001}{10.9091}{rm}{51}\RG{55.808}{0.001}{10.9091}{mi}{116}\RG{59.747}{3.960}{7.9701}{rm}{51}\RG{64.480}{0.001}{10.9091}{rm}{47}\RG{69.934}{0.001}{10.9091}{rm}{50}\RG{81.400}{0.001}{10.9091}{rm}{61}\RG{95.898}{0.001}{10.9091}{sy}{0}\RG{104.383}{0.001}{10.9091}{mi}{116}\RG{108.322}{3.960}{7.9701}{rm}{51}\RG{113.055}{0.001}{10.9091}{rm}{47}\RG{118.509}{0.001}{10.9091}{rm}{50}%
}\hskip.1em\allowbreak{} ，即\allowbreak{}\hskip.1em \RB{56.143}{8.072}{0.739}{%
\RG{0.000}{0.001}{10.9091}{mi}{116}\RG{9.951}{0.001}{10.9091}{rm}{61}\RG{24.448}{0.001}{10.9091}{mi}{100}\RG{36.137}{0.001}{10.9091}{rm}{61}\RG{50.625}{0.001}{10.9091}{rm}{48}%
}\hskip.1em\allowbreak{} 。由Cayley–Hamilton定理，\allowbreak{}\hskip.1em \RB{87.746}{9.755}{2.140}{%
\RG{0.000}{-0.004}{10.9091}{mi}{67}\RG{8.577}{3.955}{7.9701}{rm}{50}\RG{15.734}{-0.004}{10.9091}{sy}{0}\RG{26.643}{-0.004}{10.9091}{mi}{116}\RG{30.581}{-0.004}{10.9091}{mi}{67}\RG{41.584}{-0.004}{10.9091}{rm}{43}\RG{52.493}{-0.004}{10.9091}{mi}{100}\RG{58.166}{-0.004}{10.9091}{mi}{73}\RG{62.967}{-1.640}{7.9701}{rm}{50}\RG{70.730}{-0.004}{10.9091}{rm}{61}\RG{82.228}{-0.004}{10.9091}{rm}{48}%
}\hskip.1em\allowbreak{} ，故\allowbreak{}\hskip.1em \RB{33.357}{9.755}{0.744}{%
\RG{0.000}{-0.004}{10.9091}{mi}{67}\RG{8.578}{3.955}{7.9701}{rm}{50}\RG{16.340}{-0.004}{10.9091}{rm}{61}\RG{27.838}{-0.004}{10.9091}{rm}{48}%
}\hskip.1em\allowbreak{} 。从而\allowbreak{}\hskip.1em \RB{95.243}{9.755}{0.744}{%
\RG{0.000}{-0.004}{10.9091}{mi}{65}\RG{8.182}{-0.004}{10.9091}{mi}{66}\RG{19.428}{-0.004}{10.9091}{sy}{0}\RG{30.337}{-0.004}{10.9091}{mi}{66}\RG{39.152}{-0.004}{10.9091}{mi}{65}\RG{50.371}{-0.004}{10.9091}{rm}{61}\RG{61.886}{-0.004}{10.9091}{mi}{67}\RG{70.464}{3.955}{7.9701}{rm}{50}\RG{78.227}{-0.004}{10.9091}{rm}{61}\RG{89.725}{-0.004}{10.9091}{rm}{48}%
}\hskip.1em\allowbreak{} 。（2）令\allowbreak{}\hskip.1em \RB{14.722}{7.922}{3.758}{%
\RG{0.000}{0.004}{10.9091}{mi}{69}\RG{8.053}{-1.632}{7.9701}{mi}{105}\RG{10.930}{-1.632}{7.9701}{mi}{106}%
}\hskip.1em\allowbreak{} 表示\allowbreak{}\hskip.1em \RB{6.729}{5.326}{0.616}{%
\RG{0.000}{0.004}{10.9091}{mi}{110}%
}\hskip.1em\allowbreak{} 阶基础矩阵。取

\[
\RB{166.247}{8.682}{3.227}{%
\RG{0.000}{0.000}{10.9091}{mi}{65}\RG{11.212}{0.000}{10.9091}{rm}{61}\RG{22.727}{0.000}{10.9091}{mi}{69}\RG{39.747}{0.000}{10.9091}{mi}{59}\RG{66.398}{0.000}{10.9091}{mi}{66}\RG{78.267}{0.000}{10.9091}{rm}{61}\RG{89.782}{0.000}{10.9091}{mi}{69}\RG{128.620}{0.000}{10.9091}{rm}{40}\RG{132.862}{0.000}{10.9091}{mi}{110}\RG{142.440}{0.000}{10.9091}{sy}{21}\RG{153.955}{0.000}{10.9091}{rm}{51}\RG{159.410}{0.000}{10.9091}{rm}{41}\RG{163.652}{0.000}{10.9091}{mi}{58}\RG{30.780}{-1.636}{7.9701}{rm}{50}\RG{35.012}{-1.636}{7.9701}{rm}{51}\RG{97.835}{-1.636}{7.9701}{rm}{49}\RG{102.067}{-1.636}{7.9701}{rm}{50}%
}
\]
则\allowbreak{}\hskip.1em \RB{92.610}{9.755}{0.744}{%
\RG{0.000}{-0.004}{10.9091}{mi}{65}\RG{8.182}{3.955}{7.9701}{rm}{50}\RG{15.944}{-0.004}{10.9091}{rm}{61}\RG{27.459}{-0.004}{10.9091}{mi}{66}\RG{36.281}{3.955}{7.9701}{rm}{50}\RG{44.044}{-0.004}{10.9091}{rm}{61}\RG{55.559}{-0.004}{10.9091}{mi}{65}\RG{63.741}{-0.004}{10.9091}{mi}{66}\RG{75.594}{-0.004}{10.9091}{rm}{61}\RG{87.092}{-0.004}{10.9091}{rm}{48}%
}\hskip.1em\allowbreak{} ，故\allowbreak{}\hskip.1em \RB{95.640}{9.755}{1.409}{%
\RG{0.000}{-0.004}{10.9091}{mi}{65}\RG{8.181}{3.955}{7.9701}{rm}{50}\RG{15.338}{-0.004}{10.9091}{sy}{0}\RG{26.247}{-0.004}{10.9091}{rm}{50}\RG{31.702}{-0.004}{10.9091}{mi}{65}\RG{39.884}{-0.004}{10.9091}{mi}{66}\RG{51.130}{-0.004}{10.9091}{rm}{43}\RG{62.040}{-0.004}{10.9091}{mi}{66}\RG{70.862}{3.955}{7.9701}{rm}{50}\RG{78.624}{-0.004}{10.9091}{rm}{61}\RG{90.122}{-0.004}{10.9091}{rm}{48}%
}\hskip.1em\allowbreak{} ；但

\[
\RB{102.706}{8.313}{2.854}{%
\RG{0.000}{0.002}{10.9091}{mi}{66}\RG{8.815}{0.002}{10.9091}{mi}{65}\RG{20.034}{0.002}{10.9091}{rm}{61}\RG{31.549}{0.002}{10.9091}{mi}{69}\RG{51.599}{0.002}{10.9091}{sy}{54}\RG{51.599}{0.002}{10.9091}{rm}{61}\RG{63.098}{0.002}{10.9091}{rm}{48}\RG{71.585}{0.002}{10.9091}{rm}{61}\RG{83.115}{0.002}{10.9091}{mi}{65}\RG{91.297}{0.002}{10.9091}{mi}{66}\RG{100.111}{0.002}{10.9091}{mi}{58}\RG{39.602}{-1.634}{7.9701}{rm}{49}\RG{43.834}{-1.634}{7.9701}{rm}{51}%
}
\]
这给出了任意\allowbreak{}\hskip.1em \RB{26.580}{7.764}{1.997}{%
\RG{0.000}{-0.002}{10.9091}{mi}{110}\RG{9.578}{-0.002}{10.9091}{sy}{21}\RG{21.094}{-0.002}{10.9091}{rm}{51}%
}\hskip.1em\allowbreak{} 的反例。小结：把原式化为\allowbreak{}\hskip.1em \RB{55.640}{9.755}{3.231}{%
\RG{0.000}{-0.004}{10.9091}{mi}{67}\RG{8.577}{3.955}{7.9701}{rm}{50}\RG{16.779}{-0.004}{10.9091}{rm}{61}\RG{28.725}{-0.004}{10.9091}{rm}{91}\RG{31.764}{-0.004}{10.9091}{mi}{67}\RG{39.717}{-0.004}{10.9091}{mi}{59}\RG{44.561}{-0.004}{10.9091}{mi}{66}\RG{53.405}{-0.004}{10.9091}{rm}{93}%
}\hskip.1em\allowbreak{} ，再通过迹与二阶Cayley–Hamilton定理推出\allowbreak{}\hskip.1em \RB{34.243}{9.755}{0.744}{%
\RG{0.000}{-0.004}{10.9091}{mi}{67}\RG{8.577}{3.955}{7.9701}{rm}{50}\RG{16.779}{-0.004}{10.9091}{rm}{61}\RG{28.725}{-0.004}{10.9091}{rm}{48}%
}\hskip.1em\allowbreak{} ；高阶情形可由基础矩阵构造反例。

\[
\RB{278.451}{7.726}{2.294}{%
\RT{0.000}{-0.001}{7.9701}{题}\RT{8.966}{-0.001}{7.9701}{目}\RT{17.933}{-0.001}{7.9701}{来}\RT{26.899}{-0.001}{7.9701}{源}\RT{35.866}{-0.001}{7.9701}{：}\RT{44.832}{-0.001}{7.9701}{谢}\RT{53.799}{-0.001}{7.9701}{启}\RT{62.765}{-0.001}{7.9701}{鸿}\RT{71.732}{-0.001}{7.9701}{《}\RT{80.698}{-0.001}{7.9701}{复}\RT{89.665}{-0.001}{7.9701}{旦}\RT{98.631}{-0.001}{7.9701}{大}\RT{107.598}{-0.001}{7.9701}{学}\RT{116.564}{-0.001}{7.9701}{高}\RT{125.531}{-0.001}{7.9701}{等}\RT{134.497}{-0.001}{7.9701}{代}\RT{143.464}{-0.001}{7.9701}{数}\RT{152.430}{-0.001}{7.9701}{习}\RT{161.397}{-0.001}{7.9701}{题}\RT{170.363}{-0.001}{7.9701}{集}\RT{179.329}{-0.001}{7.9701}{（}\RT{188.296}{-0.001}{7.9701}{2}\RT{192.779}{-0.001}{7.9701}{0}\RT{197.262}{-0.001}{7.9701}{2}\RT{201.745}{-0.001}{7.9701}{3}\RT{209.214}{-0.001}{7.9701}{版}\RT{218.180}{-0.001}{7.9701}{）}\RT{222.672}{-0.001}{7.9701}{》}\RT{227.156}{-0.001}{7.9701}{，}\RT{236.122}{-0.001}{7.9701}{第}\RT{248.066}{-0.001}{7.9701}{7}\RT{252.549}{-0.001}{7.9701}{4}\RT{260.018}{-0.001}{7.9701}{页}\RT{268.985}{-0.001}{7.9701}{。}%
}
\]
% Original page 5.
\section*{2019S05}
设\allowbreak{}\hskip.1em \RB{8.366}{8.316}{0.500}{%
\RG{0.000}{0.005}{10.9091}{mi}{65}%
}\hskip.1em\allowbreak{} 是数域\allowbreak{}\hskip.1em \RB{10.198}{7.956}{0.500}{%
\RG{0.000}{0.005}{10.9091}{mi}{75}%
}\hskip.1em\allowbreak{} 上的\allowbreak{}\hskip.1em \RB{6.729}{5.327}{0.615}{%
\RG{0.000}{0.005}{10.9091}{mi}{110}%
}\hskip.1em\allowbreak{} 阶方阵，或具有相同行列分块方式的分块矩阵。（1）证明以下三种变换都是相似变换：（1.1）对换第\allowbreak{}\hskip.1em \RB{3.696}{7.711}{0.620}{%
\RG{0.000}{0.000}{10.9091}{mi}{105}%
}\hskip.1em\allowbreak{} 行与第\allowbreak{}\hskip.1em \RB{4.831}{7.711}{2.736}{%
\RG{0.000}{0.000}{10.9091}{mi}{106}%
}\hskip.1em\allowbreak{} 行，再对换第\allowbreak{}\hskip.1em \RB{3.696}{7.711}{0.620}{%
\RG{0.000}{0.000}{10.9091}{mi}{105}%
}\hskip.1em\allowbreak{} 列与第\allowbreak{}\hskip.1em \RB{4.831}{7.711}{2.736}{%
\RG{0.000}{0.000}{10.9091}{mi}{106}%
}\hskip.1em\allowbreak{} 列；（1.2）第\allowbreak{}\hskip.1em \RB{3.696}{7.709}{0.622}{%
\RG{0.000}{-0.002}{10.9091}{mi}{105}%
}\hskip.1em\allowbreak{} 行乘以非零常数\allowbreak{}\hskip.1em \RB{28.252}{7.949}{0.938}{%
\RG{0.000}{-0.002}{10.9091}{mi}{99}\RG{7.751}{-0.002}{10.9091}{sy}{50}\RG{18.054}{-0.002}{10.9091}{mi}{75}%
}\hskip.1em\allowbreak{} ，第\allowbreak{}\hskip.1em \RB{3.696}{7.709}{0.622}{%
\RG{0.000}{-0.002}{10.9091}{mi}{105}%
}\hskip.1em\allowbreak{} 列乘以\allowbreak{}\hskip.1em \RB{15.370}{9.757}{0.622}{%
\RG{0.000}{-0.002}{10.9091}{mi}{99}\RG{4.721}{3.957}{7.9701}{sy}{0}\RG{11.307}{3.957}{7.9701}{rm}{49}%
}\hskip.1em\allowbreak{} ；（1.3）第\allowbreak{}\hskip.1em \RB{3.696}{7.716}{0.615}{%
\RG{0.000}{0.005}{10.9091}{mi}{105}%
}\hskip.1em\allowbreak{} 行乘以\allowbreak{}\hskip.1em \RB{28.252}{7.956}{0.931}{%
\RG{0.000}{0.005}{10.9091}{mi}{99}\RG{7.751}{0.005}{10.9091}{sy}{50}\RG{18.054}{0.005}{10.9091}{mi}{75}%
}\hskip.1em\allowbreak{} 加到第\allowbreak{}\hskip.1em \RB{4.831}{7.716}{2.731}{%
\RG{0.000}{0.005}{10.9091}{mi}{106}%
}\hskip.1em\allowbreak{} 行上，第\allowbreak{}\hskip.1em \RB{4.831}{7.716}{2.731}{%
\RG{0.000}{0.005}{10.9091}{mi}{106}%
}\hskip.1em\allowbreak{} 列乘以\allowbreak{}\hskip.1em \RB{13.676}{5.327}{0.615}{%
\RG{0.000}{0.005}{10.9091}{sy}{0}\RG{8.485}{0.005}{10.9091}{mi}{99}%
}\hskip.1em\allowbreak{} 加到第\allowbreak{}\hskip.1em \RB{3.696}{7.716}{0.615}{%
\RG{0.000}{0.005}{10.9091}{mi}{105}%
}\hskip.1em\allowbreak{} 列上。（2）证明任一相似变换都是有限次上述相似初等变换的复合。（3）证明相应的三种相似分块初等变换也是相似变换：同时对换第\allowbreak{}\hskip.1em \RB{13.416}{7.712}{2.735}{%
\RG{0.000}{0.001}{10.9091}{mi}{105}\RG{3.753}{0.001}{10.9091}{mi}{59}\RG{8.586}{0.001}{10.9091}{mi}{106}%
}\hskip.1em\allowbreak{} 分块行和相应分块列；第\allowbreak{}\hskip.1em \RB{3.696}{7.707}{0.624}{%
\RG{0.000}{-0.004}{10.9091}{mi}{105}%
}\hskip.1em\allowbreak{} 分块行左乘非异阵\allowbreak{}\hskip.1em \RB{11.889}{7.947}{0.504}{%
\RG{0.000}{-0.004}{10.9091}{mi}{77}%
}\hskip.1em\allowbreak{} ，第\allowbreak{}\hskip.1em \RB{3.696}{7.707}{0.624}{%
\RG{0.000}{-0.004}{10.9091}{mi}{105}%
}\hskip.1em\allowbreak{} 分块列右乘\allowbreak{}\hskip.1em \RB{22.422}{9.755}{0.504}{%
\RG{0.000}{-0.004}{10.9091}{mi}{77}\RG{11.772}{3.955}{7.9701}{sy}{0}\RG{18.359}{3.955}{7.9701}{rm}{49}%
}\hskip.1em\allowbreak{} ；第\allowbreak{}\hskip.1em \RB{3.696}{7.707}{0.624}{%
\RG{0.000}{-0.004}{10.9091}{mi}{105}%
}\hskip.1em\allowbreak{} 分块行左乘矩阵\allowbreak{}\hskip.1em \RB{11.889}{7.947}{0.504}{%
\RG{0.000}{-0.004}{10.9091}{mi}{77}%
}\hskip.1em\allowbreak{} 加到第\allowbreak{}\hskip.1em \RB{4.831}{7.707}{2.740}{%
\RG{0.000}{-0.004}{10.9091}{mi}{106}%
}\hskip.1em\allowbreak{} 分块行，再将第\allowbreak{}\hskip.1em \RB{4.831}{7.712}{2.735}{%
\RG{0.000}{0.001}{10.9091}{mi}{106}%
}\hskip.1em\allowbreak{} 分块列右乘\allowbreak{}\hskip.1em \RB{20.373}{7.952}{0.500}{%
\RG{0.000}{0.001}{10.9091}{sy}{0}\RG{8.484}{0.001}{10.9091}{mi}{77}%
}\hskip.1em\allowbreak{} 加到第\allowbreak{}\hskip.1em \RB{3.696}{7.712}{0.619}{%
\RG{0.000}{0.001}{10.9091}{mi}{105}%
}\hskip.1em\allowbreak{} 分块列。

证明：（1）记\allowbreak{}\hskip.1em \RB{13.673}{7.949}{3.764}{%
\RG{0.000}{-0.002}{10.9091}{mi}{80}\RG{7.004}{-1.638}{7.9701}{mi}{105}\RG{9.881}{-1.638}{7.9701}{mi}{106}%
}\hskip.1em\allowbreak{} 为交换第\allowbreak{}\hskip.1em \RB{13.416}{7.709}{2.738}{%
\RG{0.000}{-0.002}{10.9091}{mi}{105}\RG{3.753}{-0.002}{10.9091}{mi}{59}\RG{8.586}{-0.002}{10.9091}{mi}{106}%
}\hskip.1em\allowbreak{} 个坐标的置换阵，\allowbreak{}\hskip.1em \RB{108.586}{8.680}{3.229}{%
\RG{0.000}{-0.002}{10.9091}{mi}{68}\RG{9.032}{-1.638}{7.9701}{mi}{105}\RG{12.414}{-0.002}{10.9091}{rm}{40}\RG{16.656}{-0.002}{10.9091}{mi}{99}\RG{21.377}{-0.002}{10.9091}{rm}{41}\RG{28.632}{-0.002}{10.9091}{rm}{61}\RG{40.165}{-0.002}{10.9091}{mi}{73}\RG{44.960}{-1.638}{7.9701}{mi}{110}\RG{52.697}{-0.002}{10.9091}{rm}{43}\RG{63.268}{-0.002}{10.9091}{rm}{40}\RG{67.526}{-0.002}{10.9091}{mi}{99}\RG{74.347}{-0.002}{10.9091}{sy}{0}\RG{84.932}{-0.002}{10.9091}{rm}{49}\RG{90.386}{-0.002}{10.9091}{rm}{41}\RG{94.629}{-0.002}{10.9091}{mi}{69}\RG{102.682}{-1.638}{7.9701}{mi}{105}\RG{105.559}{-1.638}{7.9701}{mi}{105}%
}\hskip.1em\allowbreak{} ，以及\allowbreak{}\hskip.1em \RB{84.184}{8.680}{3.764}{%
\RG{0.000}{-0.002}{10.9091}{mi}{84}\RG{6.375}{-1.638}{7.9701}{mi}{106}\RG{10.249}{-1.638}{7.9701}{mi}{105}\RG{13.641}{-0.002}{10.9091}{rm}{40}\RG{17.883}{-0.002}{10.9091}{mi}{99}\RG{22.604}{-0.002}{10.9091}{rm}{41}\RG{29.859}{-0.002}{10.9091}{rm}{61}\RG{41.392}{-0.002}{10.9091}{mi}{73}\RG{46.187}{-1.638}{7.9701}{mi}{110}\RG{53.924}{-0.002}{10.9091}{rm}{43}\RG{64.510}{-0.002}{10.9091}{mi}{99}\RG{69.222}{-0.002}{10.9091}{mi}{69}\RG{77.284}{-1.638}{7.9701}{mi}{106}\RG{81.157}{-1.638}{7.9701}{mi}{105}%
}\hskip.1em\allowbreak{}（\allowbreak{}\hskip.1em \RB{23.134}{8.314}{2.853}{%
\RG{0.000}{0.003}{10.9091}{mi}{105}\RG{6.788}{0.003}{10.9091}{sy}{54}\RG{6.788}{0.003}{10.9091}{rm}{61}\RG{18.303}{0.003}{10.9091}{mi}{106}%
}\hskip.1em\allowbreak{} ）。它们均可逆，且

\[
\RB{289.518}{10.390}{5.380}{%
\RG{8.519}{4.590}{7.9701}{sy}{0}\RG{15.106}{4.590}{7.9701}{rm}{49}\RG{100.939}{4.507}{7.9701}{sy}{0}\RG{107.525}{4.507}{7.9701}{rm}{49}\RG{148.180}{4.507}{7.9701}{sy}{0}\RG{154.766}{4.507}{7.9701}{rm}{49}\RG{217.255}{4.507}{7.9701}{sy}{0}\RG{223.841}{4.507}{7.9701}{rm}{49}\RG{0.000}{0.002}{10.9091}{mi}{80}\RG{22.868}{0.002}{10.9091}{rm}{61}\RG{34.383}{0.002}{10.9091}{mi}{80}\RG{48.652}{0.002}{10.9091}{mi}{59}\RG{75.304}{0.002}{10.9091}{mi}{68}\RG{87.733}{0.002}{10.9091}{rm}{40}\RG{91.975}{0.002}{10.9091}{mi}{99}\RG{96.696}{0.002}{10.9091}{rm}{41}\RG{115.288}{0.002}{10.9091}{rm}{61}\RG{126.803}{0.002}{10.9091}{mi}{68}\RG{139.216}{0.002}{10.9091}{rm}{40}\RG{143.459}{0.002}{10.9091}{mi}{99}\RG{159.498}{0.002}{10.9091}{rm}{41}\RG{163.741}{0.002}{10.9091}{mi}{59}\RG{190.392}{0.002}{10.9091}{mi}{84}\RG{204.049}{0.002}{10.9091}{rm}{40}\RG{208.291}{0.002}{10.9091}{mi}{99}\RG{213.012}{0.002}{10.9091}{rm}{41}\RG{231.604}{0.002}{10.9091}{rm}{61}\RG{243.119}{0.002}{10.9091}{mi}{73}\RG{255.975}{0.002}{10.9091}{sy}{0}\RG{266.884}{0.002}{10.9091}{mi}{99}\RG{271.597}{0.002}{10.9091}{mi}{69}\RG{286.923}{0.002}{10.9091}{mi}{58}\RG{41.387}{-1.634}{7.9701}{mi}{105}\RG{44.265}{-1.634}{7.9701}{mi}{106}\RG{84.352}{-1.634}{7.9701}{mi}{105}\RG{135.835}{-1.634}{7.9701}{mi}{105}\RG{196.783}{-1.634}{7.9701}{mi}{106}\RG{200.656}{-1.634}{7.9701}{mi}{105}\RG{247.914}{-1.634}{7.9701}{mi}{110}\RG{279.658}{-1.634}{7.9701}{mi}{106}\RG{283.532}{-1.634}{7.9701}{mi}{105}\RG{7.004}{-3.254}{7.9701}{mi}{105}\RG{9.881}{-3.254}{7.9701}{mi}{106}%
}
\]
于是三类操作分别表示为

\[
\RB{107.606}{29.798}{21.930}{%
\RG{62.789}{23.998}{7.9701}{sy}{0}\RG{69.376}{23.998}{7.9701}{rm}{49}\RG{0.000}{19.410}{10.9091}{mi}{65}\RG{11.212}{19.410}{10.9091}{sy}{55}\RG{11.212}{19.410}{10.9091}{sy}{0}\RG{17.878}{19.410}{10.9091}{sy}{33}\RG{31.818}{19.410}{10.9091}{mi}{80}\RG{46.087}{19.410}{10.9091}{mi}{65}\RG{54.269}{19.410}{10.9091}{mi}{80}\RG{74.108}{19.410}{10.9091}{mi}{59}\RG{38.822}{17.774}{7.9701}{mi}{105}\RG{41.700}{17.774}{7.9701}{mi}{106}\RG{61.274}{16.154}{7.9701}{mi}{105}\RG{64.151}{16.154}{7.9701}{mi}{106}\RG{91.239}{4.501}{7.9701}{sy}{0}\RG{97.825}{4.501}{7.9701}{rm}{49}\RG{0.000}{-0.004}{10.9091}{mi}{65}\RG{11.212}{-0.004}{10.9091}{sy}{55}\RG{11.212}{-0.004}{10.9091}{sy}{0}\RG{17.878}{-0.004}{10.9091}{sy}{33}\RG{31.818}{-0.004}{10.9091}{mi}{68}\RG{44.231}{-0.004}{10.9091}{rm}{40}\RG{48.474}{-0.004}{10.9091}{mi}{99}\RG{53.195}{-0.004}{10.9091}{rm}{41}\RG{57.438}{-0.004}{10.9091}{mi}{65}\RG{65.619}{-0.004}{10.9091}{mi}{68}\RG{78.033}{-0.004}{10.9091}{rm}{40}\RG{82.276}{-0.004}{10.9091}{mi}{99}\RG{86.997}{-0.004}{10.9091}{rm}{41}\RG{102.558}{-0.004}{10.9091}{mi}{59}\RG{40.850}{-1.640}{7.9701}{mi}{105}\RG{74.652}{-1.640}{7.9701}{mi}{105}\RG{93.693}{-13.663}{7.9701}{sy}{0}\RG{100.280}{-13.663}{7.9701}{rm}{49}\RG{0.000}{-18.168}{10.9091}{mi}{65}\RG{11.212}{-18.168}{10.9091}{sy}{55}\RG{11.212}{-18.168}{10.9091}{sy}{0}\RG{17.878}{-18.168}{10.9091}{sy}{33}\RG{31.818}{-18.168}{10.9091}{mi}{84}\RG{45.459}{-18.168}{10.9091}{rm}{40}\RG{49.701}{-18.168}{10.9091}{mi}{99}\RG{54.422}{-18.168}{10.9091}{rm}{41}\RG{58.665}{-18.168}{10.9091}{mi}{65}\RG{66.846}{-18.168}{10.9091}{mi}{84}\RG{80.488}{-18.168}{10.9091}{rm}{40}\RG{84.730}{-18.168}{10.9091}{mi}{99}\RG{89.451}{-18.168}{10.9091}{rm}{41}\RG{105.012}{-18.168}{10.9091}{mi}{58}\RG{38.193}{-19.804}{7.9701}{mi}{106}\RG{42.067}{-19.804}{7.9701}{mi}{105}\RG{73.222}{-19.804}{7.9701}{mi}{106}\RG{77.095}{-19.804}{7.9701}{mi}{105}%
}
\]
其中第三式左乘\allowbreak{}\hskip.1em \RB{26.257}{8.683}{3.761}{%
\RG{0.000}{0.001}{10.9091}{mi}{84}\RG{6.375}{-1.635}{7.9701}{mi}{106}\RG{10.249}{-1.635}{7.9701}{mi}{105}\RG{13.640}{0.001}{10.9091}{rm}{40}\RG{17.883}{0.001}{10.9091}{mi}{99}\RG{22.604}{0.001}{10.9091}{rm}{41}%
}\hskip.1em\allowbreak{} 是将第\allowbreak{}\hskip.1em \RB{3.696}{7.712}{0.619}{%
\RG{0.000}{0.001}{10.9091}{mi}{105}%
}\hskip.1em\allowbreak{} 行的\allowbreak{}\hskip.1em \RB{5.191}{5.323}{0.619}{%
\RG{0.000}{0.001}{10.9091}{mi}{99}%
}\hskip.1em\allowbreak{} 倍加到第\allowbreak{}\hskip.1em \RB{4.831}{7.712}{2.735}{%
\RG{0.000}{0.001}{10.9091}{mi}{106}%
}\hskip.1em\allowbreak{} 行；右乘其逆阵则是将第\allowbreak{}\hskip.1em \RB{4.831}{7.712}{2.735}{%
\RG{0.000}{0.001}{10.9091}{mi}{106}%
}\hskip.1em\allowbreak{} 列的\allowbreak{}\hskip.1em \RB{13.675}{5.323}{0.619}{%
\RG{0.000}{0.001}{10.9091}{sy}{0}\RG{8.484}{0.001}{10.9091}{mi}{99}%
}\hskip.1em\allowbreak{} 倍加到第\allowbreak{}\hskip.1em \RB{3.696}{7.707}{0.624}{%
\RG{0.000}{-0.004}{10.9091}{mi}{105}%
}\hskip.1em\allowbreak{} 列。故三类变换都是相似变换。（2）任取非异阵\allowbreak{}\hskip.1em \RB{18.842}{7.955}{0.932}{%
\RG{0.000}{0.004}{10.9091}{mi}{80}\RG{11.982}{0.004}{10.9091}{sy}{50}%
}\hskip.1em\allowbreak{} GL\allowbreak{}\hskip.1em \RB{23.578}{8.686}{3.223}{%
\RG{0.000}{-1.632}{7.9701}{mi}{110}\RG{5.637}{0.004}{10.9091}{rm}{40}\RG{9.879}{0.004}{10.9091}{mi}{75}\RG{19.925}{0.004}{10.9091}{rm}{41}%
}\hskip.1em\allowbreak{} 。通过有限次初等行变换可将\allowbreak{}\hskip.1em \RB{8.726}{7.955}{0.500}{%
\RG{0.000}{0.004}{10.9091}{mi}{80}%
}\hskip.1em\allowbreak{} 化为\allowbreak{}\hskip.1em \RB{10.174}{7.955}{2.212}{%
\RG{0.000}{0.004}{10.9091}{mi}{73}\RG{4.796}{-1.632}{7.9701}{mi}{110}%
}\hskip.1em\allowbreak{} ，所以\allowbreak{}\hskip.1em \RB{8.726}{7.955}{0.500}{%
\RG{0.000}{0.004}{10.9091}{mi}{80}%
}\hskip.1em\allowbreak{} 可写为有限个第一、二、三类初等阵的乘积\allowbreak{}\hskip.1em \RB{64.011}{7.950}{2.217}{%
\RG{0.000}{-0.001}{10.9091}{mi}{80}\RG{11.549}{-0.001}{10.9091}{rm}{61}\RG{23.064}{-0.001}{10.9091}{mi}{69}\RG{31.117}{-1.637}{7.9701}{mi}{115}\RG{37.349}{-0.001}{10.9091}{sy}{1}\RG{42.182}{-0.001}{10.9091}{sy}{1}\RG{47.026}{-0.001}{10.9091}{sy}{1}\RG{51.895}{-0.001}{10.9091}{mi}{69}\RG{59.948}{-1.637}{7.9701}{rm}{49}%
}\hskip.1em\allowbreak{} 。因此

\[
\RB{213.841}{13.044}{7.582}{%
\RG{63.552}{12.107}{10.9092}{ex}{16}\RG{184.731}{12.107}{10.9092}{ex}{17}\RG{97.398}{8.834}{10.9092}{ex}{0}\RG{143.366}{8.834}{10.9092}{ex}{1}\RG{132.047}{4.586}{7.9701}{sy}{0}\RG{138.634}{4.586}{7.9701}{rm}{49}\RG{157.048}{4.586}{7.9701}{sy}{0}\RG{163.634}{4.586}{7.9701}{rm}{49}\RG{25.220}{4.502}{7.9701}{sy}{0}\RG{31.807}{4.502}{7.9701}{rm}{49}\RG{199.928}{4.502}{7.9701}{sy}{0}\RG{206.515}{4.502}{7.9701}{rm}{49}\RG{0.000}{-0.002}{10.9091}{mi}{80}\RG{8.509}{-0.002}{10.9091}{mi}{65}\RG{16.691}{-0.002}{10.9091}{mi}{80}\RG{39.569}{-0.002}{10.9091}{rm}{61}\RG{51.084}{-0.002}{10.9091}{mi}{69}\RG{70.068}{-0.002}{10.9091}{sy}{1}\RG{74.900}{-0.002}{10.9091}{sy}{1}\RG{79.744}{-0.002}{10.9091}{sy}{1}\RG{84.613}{-0.002}{10.9091}{mi}{69}\RG{102.398}{-0.002}{10.9091}{mi}{69}\RG{115.184}{-0.002}{10.9091}{mi}{65}\RG{123.366}{-0.002}{10.9091}{mi}{69}\RG{148.366}{-0.002}{10.9091}{mi}{69}\RG{170.185}{-0.002}{10.9091}{sy}{1}\RG{175.018}{-0.002}{10.9091}{sy}{1}\RG{179.862}{-0.002}{10.9091}{sy}{1}\RG{191.246}{-0.002}{10.9091}{mi}{69}\RG{211.247}{-0.002}{10.9091}{mi}{58}\RG{59.137}{-1.638}{7.9701}{mi}{115}\RG{92.666}{-1.638}{7.9701}{rm}{50}\RG{110.451}{-1.638}{7.9701}{rm}{49}\RG{199.299}{-2.699}{7.9701}{mi}{115}\RG{131.418}{-3.126}{7.9701}{rm}{49}\RG{156.419}{-3.126}{7.9701}{rm}{50}%
}
\]
由（1），右式就是有限次相似初等变换的复合。（3）设分块大小为\allowbreak{}\hskip.1em \RB{46.380}{5.317}{3.687}{%
\RG{0.000}{-0.005}{10.9091}{mi}{110}\RG{6.548}{-1.641}{7.9701}{rm}{49}\RG{11.280}{-0.005}{10.9091}{mi}{59}\RG{16.113}{-0.005}{10.9091}{mi}{58}\RG{20.957}{-0.005}{10.9091}{mi}{58}\RG{25.789}{-0.005}{10.9091}{mi}{58}\RG{30.633}{-0.005}{10.9091}{mi}{59}\RG{35.477}{-0.005}{10.9091}{mi}{110}\RG{42.070}{-1.641}{7.9701}{mi}{113}%
}\hskip.1em\allowbreak{} 。同时对换第\allowbreak{}\hskip.1em \RB{13.416}{7.706}{2.741}{%
\RG{0.000}{-0.005}{10.9091}{mi}{105}\RG{3.753}{-0.005}{10.9091}{mi}{59}\RG{8.586}{-0.005}{10.9091}{mi}{106}%
}\hskip.1em\allowbreak{} 个坐标块，可用相应的坐标置换阵\allowbreak{}\hskip.1em \RB{8.726}{7.946}{0.505}{%
\RG{0.000}{-0.005}{10.9091}{mi}{80}%
}\hskip.1em\allowbreak{} 作\allowbreak{}\hskip.1em \RB{61.021}{9.759}{0.620}{%
\RG{0.000}{0.000}{10.9091}{mi}{65}\RG{11.212}{0.000}{10.9091}{sy}{55}\RG{11.212}{0.000}{10.9091}{sy}{33}\RG{25.151}{0.000}{10.9091}{mi}{80}\RG{33.660}{0.000}{10.9091}{mi}{65}\RG{41.842}{0.000}{10.9091}{mi}{80}\RG{50.371}{3.959}{7.9701}{sy}{0}\RG{56.958}{3.959}{7.9701}{rm}{49}%
}\hskip.1em\allowbreak{} ；若块的大小不同，交换后同时按新的顺序标记分块。

对于第二类分块变换，令

\[
\RB{199.415}{9.290}{4.296}{%
\RG{0.000}{-0.003}{10.9091}{mi}{68}\RG{12.366}{-0.003}{10.9091}{rm}{61}\RT{23.881}{-0.003}{10.9091}{d}\RT{29.947}{-0.003}{10.9091}{i}\RT{32.979}{-0.003}{10.9091}{a}\RT{38.434}{-0.003}{10.9091}{g}\RG{43.888}{-0.003}{10.9091}{rm}{40}\RG{48.130}{-0.003}{10.9091}{mi}{73}\RG{62.714}{-0.003}{10.9091}{mi}{59}\RG{67.546}{-0.003}{10.9091}{mi}{58}\RG{72.390}{-0.003}{10.9091}{mi}{58}\RG{77.234}{-0.003}{10.9091}{mi}{58}\RG{82.066}{-0.003}{10.9091}{mi}{59}\RG{86.910}{-0.003}{10.9091}{mi}{73}\RG{109.959}{-0.003}{10.9091}{mi}{59}\RG{114.792}{-0.003}{10.9091}{mi}{77}\RG{125.963}{-0.003}{10.9091}{mi}{59}\RG{130.795}{-0.003}{10.9091}{mi}{73}\RG{153.659}{-0.003}{10.9091}{mi}{59}\RG{158.492}{-0.003}{10.9091}{mi}{58}\RG{163.336}{-0.003}{10.9091}{mi}{58}\RG{168.179}{-0.003}{10.9091}{mi}{58}\RG{173.012}{-0.003}{10.9091}{mi}{59}\RG{177.856}{-0.003}{10.9091}{mi}{73}\RG{192.458}{-0.003}{10.9091}{rm}{41}\RG{196.701}{-0.003}{10.9091}{mi}{59}\RG{52.926}{-1.640}{7.9701}{mi}{110}\RG{91.752}{-1.640}{7.9701}{mi}{110}\RG{135.618}{-1.640}{7.9701}{mi}{110}\RG{182.698}{-1.640}{7.9701}{mi}{110}\RG{187.836}{-2.636}{5.9776}{mi}{113}\RG{58.064}{-2.747}{5.9776}{rm}{49}\RG{96.890}{-2.855}{5.9776}{mi}{105}\RG{99.554}{-2.855}{5.9776}{sy}{0}\RG{105.310}{-2.855}{5.9776}{rm}{49}\RG{140.756}{-2.855}{5.9776}{mi}{105}\RG{143.420}{-2.855}{5.9776}{rm}{43}\RG{149.009}{-2.855}{5.9776}{rm}{49}%
}
\]
其中\allowbreak{}\hskip.1em \RB{21.663}{7.947}{0.940}{%
\RG{0.000}{-0.004}{10.9091}{mi}{77}\RG{14.803}{-0.004}{10.9091}{sy}{50}%
}\hskip.1em\allowbreak{} GL\allowbreak{}\hskip.1em \RB{26.739}{8.678}{3.416}{%
\RG{0.000}{-1.641}{7.9701}{mi}{110}\RG{5.138}{-2.856}{5.9776}{mi}{105}\RG{8.798}{-0.004}{10.9091}{rm}{40}\RG{13.040}{-0.004}{10.9091}{mi}{75}\RG{23.086}{-0.004}{10.9091}{rm}{41}%
}\hskip.1em\allowbreak{} 。则\allowbreak{}\hskip.1em \RB{62.652}{9.755}{0.624}{%
\RG{0.000}{-0.004}{10.9091}{mi}{65}\RG{11.212}{-0.004}{10.9091}{sy}{55}\RG{11.212}{-0.004}{10.9091}{sy}{33}\RG{25.151}{-0.004}{10.9091}{mi}{68}\RG{34.468}{-0.004}{10.9091}{mi}{65}\RG{42.650}{-0.004}{10.9091}{mi}{68}\RG{52.003}{3.955}{7.9701}{sy}{0}\RG{58.589}{3.955}{7.9701}{rm}{49}%
}\hskip.1em\allowbreak{} 恰好是题述的行、列分块操作。

对于第三类分块变换，取任意\allowbreak{}\hskip.1em \RB{78.921}{8.685}{4.568}{%
\RG{0.000}{0.003}{10.9091}{mi}{77}\RG{15.236}{0.003}{10.9091}{sy}{50}\RG{25.971}{0.003}{10.9091}{mi}{77}\RG{36.555}{-1.633}{7.9701}{mi}{110}\RG{41.693}{-2.848}{5.9776}{mi}{106}\RG{45.596}{-1.633}{7.9701}{sy}{2}\RG{52.182}{-1.633}{7.9701}{mi}{110}\RG{57.320}{-2.848}{5.9776}{mi}{105}\RG{60.980}{0.003}{10.9091}{rm}{40}\RG{65.224}{0.003}{10.9091}{mi}{75}\RG{75.269}{0.003}{10.9091}{rm}{41}%
}\hskip.1em\allowbreak{} ，定义分块矩阵\allowbreak{}\hskip.1em \RB{8.791}{7.954}{0.737}{%
\RG{0.000}{0.003}{10.9091}{mi}{85}%
}\hskip.1em\allowbreak{} 的\allowbreak{}\hskip.1em \RB{21.013}{8.685}{3.224}{%
\RG{0.000}{0.003}{10.9091}{rm}{40}\RG{4.242}{0.003}{10.9091}{mi}{106}\RG{8.737}{0.003}{10.9091}{mi}{59}\RG{13.580}{0.003}{10.9091}{mi}{105}\RG{17.360}{0.003}{10.9091}{rm}{41}%
}\hskip.1em\allowbreak{} 分块等于\allowbreak{}\hskip.1em \RB{11.889}{7.954}{0.500}{%
\RG{0.000}{0.003}{10.9091}{mi}{77}%
}\hskip.1em\allowbreak{} ，其余分块均为零。由于\allowbreak{}\hskip.1em \RB{23.135}{8.309}{2.859}{%
\RG{0.000}{-0.002}{10.9091}{mi}{105}\RG{6.789}{-0.002}{10.9091}{sy}{54}\RG{6.789}{-0.002}{10.9091}{rm}{61}\RG{18.304}{-0.002}{10.9091}{mi}{106}%
}\hskip.1em\allowbreak{} ，有\allowbreak{}\hskip.1em \RB{33.417}{9.757}{0.742}{%
\RG{0.000}{-0.002}{10.9091}{mi}{85}\RG{8.638}{3.957}{7.9701}{rm}{50}\RG{16.400}{-0.002}{10.9091}{rm}{61}\RG{27.898}{-0.002}{10.9091}{rm}{48}%
}\hskip.1em\allowbreak{} ，故

\[
\RB{149.410}{10.302}{2.609}{%
\RG{89.767}{4.502}{7.9701}{sy}{0}\RG{96.354}{4.502}{7.9701}{rm}{49}\RG{0.000}{-0.003}{10.9091}{mi}{69}\RG{11.712}{-0.003}{10.9091}{rm}{61}\RG{23.227}{-0.003}{10.9091}{mi}{73}\RG{36.083}{-0.003}{10.9091}{rm}{43}\RG{46.992}{-0.003}{10.9091}{mi}{85}\RG{54.411}{-0.003}{10.9091}{mi}{59}\RG{81.062}{-0.003}{10.9091}{mi}{69}\RG{104.116}{-0.003}{10.9091}{rm}{61}\RG{115.632}{-0.003}{10.9091}{mi}{73}\RG{128.488}{-0.003}{10.9091}{sy}{0}\RG{139.397}{-0.003}{10.9091}{mi}{85}\RG{146.815}{-0.003}{10.9091}{mi}{58}\RG{28.023}{-1.639}{7.9701}{mi}{110}\RG{120.427}{-1.639}{7.9701}{mi}{110}%
}
\]
左乘\allowbreak{}\hskip.1em \RB{8.832}{7.914}{0.504}{%
\RG{0.000}{-0.004}{10.9091}{mi}{69}%
}\hskip.1em\allowbreak{} 将第\allowbreak{}\hskip.1em \RB{3.696}{7.707}{0.624}{%
\RG{0.000}{-0.004}{10.9091}{mi}{105}%
}\hskip.1em\allowbreak{} 分块行左乘\allowbreak{}\hskip.1em \RB{11.889}{7.947}{0.504}{%
\RG{0.000}{-0.004}{10.9091}{mi}{77}%
}\hskip.1em\allowbreak{} 加到第\allowbreak{}\hskip.1em \RB{4.831}{7.707}{2.741}{%
\RG{0.000}{-0.004}{10.9091}{mi}{106}%
}\hskip.1em\allowbreak{} 分块行；右乘\allowbreak{}\hskip.1em \RB{19.332}{9.755}{0.504}{%
\RG{0.000}{-0.004}{10.9091}{mi}{69}\RG{8.682}{3.955}{7.9701}{sy}{0}\RG{15.269}{3.955}{7.9701}{rm}{49}%
}\hskip.1em\allowbreak{} 将第\allowbreak{}\hskip.1em \RB{4.831}{7.707}{2.741}{%
\RG{0.000}{-0.004}{10.9091}{mi}{106}%
}\hskip.1em\allowbreak{} 分块列右乘\allowbreak{}\hskip.1em \RB{20.374}{7.947}{0.504}{%
\RG{0.000}{-0.004}{10.9091}{sy}{0}\RG{8.485}{-0.004}{10.9091}{mi}{77}%
}\hskip.1em\allowbreak{} 加到第\allowbreak{}\hskip.1em \RB{3.696}{7.707}{0.624}{%
\RG{0.000}{-0.004}{10.9091}{mi}{105}%
}\hskip.1em\allowbreak{} 分块列。因此三种相似分块初等变换也都是相似变换。小结：三类普通或分块变换都可统一写为\allowbreak{}\hskip.1em \RB{62.169}{9.758}{0.621}{%
\RG{0.000}{-0.001}{10.9091}{mi}{65}\RG{11.623}{-0.001}{10.9091}{sy}{55}\RG{11.623}{-0.001}{10.9091}{sy}{33}\RG{25.974}{-0.001}{10.9091}{mi}{69}\RG{34.647}{-0.001}{10.9091}{mi}{65}\RG{42.829}{-0.001}{10.9091}{mi}{69}\RG{51.519}{3.958}{7.9701}{sy}{0}\RG{58.106}{3.958}{7.9701}{rm}{49}%
}\hskip.1em\allowbreak{} ；再利用可逆阵的初等分解证明任意相似变换均由它们复合而成。

\[
\RB{278.451}{7.729}{2.291}{%
\RT{0.000}{0.002}{7.9701}{题}\RT{8.966}{0.002}{7.9701}{目}\RT{17.933}{0.002}{7.9701}{来}\RT{26.899}{0.002}{7.9701}{源}\RT{35.866}{0.002}{7.9701}{：}\RT{44.832}{0.002}{7.9701}{谢}\RT{53.799}{0.002}{7.9701}{启}\RT{62.765}{0.002}{7.9701}{鸿}\RT{71.732}{0.002}{7.9701}{《}\RT{80.698}{0.002}{7.9701}{复}\RT{89.665}{0.002}{7.9701}{旦}\RT{98.631}{0.002}{7.9701}{大}\RT{107.598}{0.002}{7.9701}{学}\RT{116.564}{0.002}{7.9701}{高}\RT{125.531}{0.002}{7.9701}{等}\RT{134.497}{0.002}{7.9701}{代}\RT{143.464}{0.002}{7.9701}{数}\RT{152.430}{0.002}{7.9701}{习}\RT{161.397}{0.002}{7.9701}{题}\RT{170.363}{0.002}{7.9701}{集}\RT{179.329}{0.002}{7.9701}{（}\RT{188.296}{0.002}{7.9701}{2}\RT{192.779}{0.002}{7.9701}{0}\RT{197.262}{0.002}{7.9701}{2}\RT{201.745}{0.002}{7.9701}{3}\RT{209.214}{0.002}{7.9701}{版}\RT{218.180}{0.002}{7.9701}{）}\RT{222.672}{0.002}{7.9701}{》}\RT{227.156}{0.002}{7.9701}{，}\RT{236.122}{0.002}{7.9701}{第}\RT{248.066}{0.002}{7.9701}{7}\RT{252.549}{0.002}{7.9701}{4}\RT{260.018}{0.002}{7.9701}{页}\RT{268.985}{0.002}{7.9701}{。}%
}
\]
% Original page 6.
\section*{2019S06}
设\allowbreak{}\hskip.1em \RB{55.675}{8.684}{3.225}{%
\RG{0.000}{0.002}{10.9091}{mi}{65}\RG{11.212}{0.002}{10.9091}{sy}{50}\RG{21.515}{0.002}{10.9091}{mi}{77}\RG{32.098}{-1.635}{7.9701}{mi}{110}\RG{37.735}{0.002}{10.9091}{rm}{40}\RG{41.977}{0.002}{10.9091}{mi}{75}\RG{52.022}{0.002}{10.9091}{rm}{41}%
}\hskip.1em\allowbreak{} 、\allowbreak{}\hskip.1em \RB{70.393}{8.684}{3.225}{%
\RG{0.000}{0.002}{10.9091}{mi}{66}\RG{11.852}{0.002}{10.9091}{sy}{50}\RG{22.155}{0.002}{10.9091}{mi}{77}\RG{32.739}{-1.635}{7.9701}{mi}{110}\RG{37.877}{-1.635}{7.9701}{sy}{2}\RG{44.463}{-1.635}{7.9701}{mi}{109}\RG{52.452}{0.002}{10.9091}{rm}{40}\RG{56.694}{0.002}{10.9091}{mi}{75}\RG{66.741}{0.002}{10.9091}{rm}{41}%
}\hskip.1em\allowbreak{} ，分块矩阵

\[
\RB{155.509}{10.302}{3.229}{%
\RG{88.268}{4.502}{7.9701}{mi}{110}\RG{93.406}{4.502}{7.9701}{sy}{0}\RG{99.993}{4.502}{7.9701}{rm}{50}\RG{126.577}{4.502}{7.9701}{mi}{110}\RG{131.715}{4.502}{7.9701}{sy}{0}\RG{138.302}{4.502}{7.9701}{rm}{49}\RG{0.000}{-0.002}{10.9091}{sy}{67}\RG{9.411}{-0.002}{10.9091}{rm}{61}\RG{20.909}{-0.002}{10.9091}{rm}{40}\RG{25.168}{-0.002}{10.9091}{mi}{66}\RG{33.983}{-0.002}{10.9091}{mi}{59}\RG{38.816}{-0.002}{10.9091}{mi}{65}\RG{46.997}{-0.002}{10.9091}{mi}{66}\RG{55.823}{-0.002}{10.9091}{mi}{59}\RG{60.656}{-0.002}{10.9091}{mi}{58}\RG{65.500}{-0.002}{10.9091}{mi}{58}\RG{70.343}{-0.002}{10.9091}{mi}{58}\RG{75.176}{-0.002}{10.9091}{mi}{59}\RG{80.020}{-0.002}{10.9091}{mi}{65}\RG{104.725}{-0.002}{10.9091}{mi}{66}\RG{113.539}{-0.002}{10.9091}{mi}{59}\RG{118.372}{-0.002}{10.9091}{mi}{65}\RG{143.034}{-0.002}{10.9091}{mi}{66}\RG{151.856}{-0.002}{10.9091}{rm}{41}%
}
\]
的秩为\allowbreak{}\hskip.1em \RB{5.256}{5.326}{0.616}{%
\RG{0.000}{0.004}{10.9091}{mi}{114}%
}\hskip.1em\allowbreak{} 。证明：存在可逆阵\allowbreak{}\hskip.1em \RB{56.013}{8.686}{3.223}{%
\RG{0.000}{0.004}{10.9091}{mi}{80}\RG{11.549}{0.004}{10.9091}{sy}{50}\RG{21.852}{0.004}{10.9091}{mi}{77}\RG{32.436}{-1.633}{7.9701}{mi}{110}\RG{38.072}{0.004}{10.9091}{rm}{40}\RG{42.315}{0.004}{10.9091}{mi}{75}\RG{52.360}{0.004}{10.9091}{rm}{41}%
}\hskip.1em\allowbreak{} ，使得

\[
\RB{219.120}{19.592}{14.125}{%
\RG{51.084}{18.655}{10.9092}{ex}{32}\RG{103.980}{18.655}{10.9092}{ex}{33}\RG{184.308}{18.655}{10.9092}{ex}{32}\RG{205.951}{18.655}{10.9092}{ex}{33}\RG{59.720}{7.280}{10.9091}{mi}{65}\RG{86.832}{7.280}{10.9091}{mi}{65}\RG{192.944}{7.280}{10.9091}{mi}{66}\RG{67.903}{5.644}{7.9701}{rm}{49}\RG{72.135}{5.644}{7.9701}{rm}{49}\RG{95.014}{5.644}{7.9701}{rm}{49}\RG{99.246}{5.644}{7.9701}{rm}{50}\RG{201.219}{5.644}{7.9701}{rm}{49}\RG{8.519}{4.505}{7.9701}{sy}{0}\RG{15.105}{4.505}{7.9701}{rm}{49}\RG{149.621}{4.505}{7.9701}{sy}{0}\RG{156.207}{4.505}{7.9701}{rm}{49}\RG{0.000}{0.000}{10.9091}{mi}{80}\RG{19.838}{0.000}{10.9091}{mi}{65}\RG{28.020}{0.000}{10.9091}{mi}{80}\RG{39.569}{0.000}{10.9091}{rm}{61}\RG{114.435}{0.000}{10.9091}{mi}{59}\RG{141.086}{0.000}{10.9091}{mi}{80}\RG{160.939}{0.000}{10.9091}{mi}{66}\RG{172.792}{0.000}{10.9091}{rm}{61}\RG{216.406}{0.000}{10.9091}{mi}{59}\RG{65.568}{-7.895}{10.9091}{rm}{48}\RG{86.832}{-7.895}{10.9091}{mi}{65}\RG{196.720}{-7.895}{10.9091}{rm}{48}\RG{95.014}{-9.531}{7.9701}{rm}{50}\RG{99.246}{-9.531}{7.9701}{rm}{50}%
}
\]
其中\allowbreak{}\hskip.1em \RB{63.561}{8.683}{3.226}{%
\RG{0.000}{0.001}{10.9091}{mi}{65}\RG{8.182}{-1.635}{7.9701}{rm}{49}\RG{12.414}{-1.635}{7.9701}{rm}{49}\RG{20.179}{0.001}{10.9091}{sy}{50}\RG{30.482}{0.001}{10.9091}{mi}{77}\RG{41.065}{-1.635}{7.9701}{mi}{114}\RG{45.620}{0.001}{10.9091}{rm}{40}\RG{49.862}{0.001}{10.9091}{mi}{75}\RG{59.908}{0.001}{10.9091}{rm}{41}%
}\hskip.1em\allowbreak{} ，\allowbreak{}\hskip.1em \RB{73.496}{8.683}{3.226}{%
\RG{0.000}{0.001}{10.9091}{mi}{66}\RG{8.275}{-1.635}{7.9701}{rm}{49}\RG{16.037}{0.001}{10.9091}{sy}{50}\RG{26.340}{0.001}{10.9091}{mi}{77}\RG{36.924}{-1.635}{7.9701}{mi}{114}\RG{40.980}{-1.635}{7.9701}{sy}{2}\RG{47.567}{-1.635}{7.9701}{mi}{109}\RG{55.555}{0.001}{10.9091}{rm}{40}\RG{59.798}{0.001}{10.9091}{mi}{75}\RG{69.844}{0.001}{10.9091}{rm}{41}%
}\hskip.1em\allowbreak{} 。

证明：令

\[
\RB{155.630}{19.439}{15.210}{%
\RG{62.134}{13.639}{7.9701}{mi}{110}\RG{67.272}{13.639}{7.9701}{sy}{0}\RG{73.859}{13.639}{7.9701}{rm}{49}\RG{62.235}{10.367}{10.9092}{ex}{88}\RG{105.361}{4.507}{7.9701}{mi}{106}\RG{147.398}{4.507}{7.9701}{mi}{110}\RG{0.000}{0.003}{10.9091}{mi}{87}\RG{14.849}{0.003}{10.9091}{rm}{61}\RT{26.364}{0.003}{10.9091}{I}\RT{30.302}{0.003}{10.9091}{m}\RG{41.207}{0.003}{10.9091}{sy}{67}\RG{50.618}{0.003}{10.9091}{rm}{61}\RT{79.911}{0.003}{10.9091}{I}\RT{83.849}{0.003}{10.9091}{m}\RG{92.936}{0.003}{10.9091}{rm}{40}\RG{97.179}{0.003}{10.9091}{mi}{65}\RG{109.743}{0.003}{10.9091}{mi}{66}\RG{118.565}{0.003}{10.9091}{rm}{41}\RG{125.838}{0.003}{10.9091}{sy}{18}\RG{137.353}{0.003}{10.9091}{mi}{75}\RG{153.035}{0.003}{10.9091}{mi}{58}\RG{62.761}{-13.084}{7.9701}{mi}{106}\RG{66.645}{-13.084}{7.9701}{rm}{61}\RG{73.228}{-13.084}{7.9701}{rm}{48}%
}
\]
由题设，dim\allowbreak{}\hskip.1em \RB{23.214}{7.951}{0.740}{%
\RG{0.000}{0.000}{10.9091}{mi}{87}\RG{14.848}{0.000}{10.9091}{rm}{61}%
}\hskip.1em\allowbreak{} rank\allowbreak{}\hskip.1em \RB{26.183}{8.191}{0.762}{%
\RG{0.000}{0.000}{10.9091}{sy}{67}\RG{9.411}{0.000}{10.9091}{rm}{61}\RG{20.926}{0.000}{10.9091}{mi}{114}%
}\hskip.1em\allowbreak{} ，并且Im\allowbreak{}\hskip.1em \RB{35.301}{7.951}{1.995}{%
\RG{0.000}{0.000}{10.9091}{mi}{66}\RG{11.852}{0.000}{10.9091}{sy}{18}\RG{23.368}{0.000}{10.9091}{mi}{87}%
}\hskip.1em\allowbreak{} 。

设\allowbreak{}\hskip.1em \RB{8.366}{8.309}{0.502}{%
\RG{0.000}{-0.002}{10.9091}{mi}{65}%
}\hskip.1em\allowbreak{} 的特征多项式为

\[
\RB{155.271}{10.307}{3.224}{%
\RG{40.989}{4.507}{7.9701}{mi}{110}\RG{87.372}{4.507}{7.9701}{mi}{110}\RG{92.511}{4.507}{7.9701}{sy}{0}\RG{99.097}{4.507}{7.9701}{rm}{49}\RG{0.000}{0.003}{10.9091}{mi}{112}\RG{5.489}{0.003}{10.9091}{rm}{40}\RG{9.731}{0.003}{10.9091}{mi}{120}\RG{15.966}{0.003}{10.9091}{rm}{41}\RG{23.221}{0.003}{10.9091}{rm}{61}\RG{34.754}{0.003}{10.9091}{mi}{120}\RG{49.049}{0.003}{10.9091}{rm}{43}\RG{59.959}{0.003}{10.9091}{mi}{99}\RG{81.137}{0.003}{10.9091}{mi}{120}\RG{106.254}{0.003}{10.9091}{rm}{43}\RG{117.163}{0.003}{10.9091}{sy}{1}\RG{121.996}{0.003}{10.9091}{sy}{1}\RG{126.839}{0.003}{10.9091}{sy}{1}\RG{132.314}{0.003}{10.9091}{rm}{43}\RG{143.223}{0.003}{10.9091}{mi}{99}\RG{152.676}{0.003}{10.9091}{mi}{58}\RG{64.680}{-1.633}{7.9701}{mi}{110}\RG{69.819}{-1.633}{7.9701}{sy}{0}\RG{76.405}{-1.633}{7.9701}{rm}{49}\RG{147.944}{-1.633}{7.9701}{rm}{48}%
}
\]
由Cayley–Hamilton定理\allowbreak{}\hskip.1em \RB{42.195}{8.684}{3.225}{%
\RG{0.000}{0.002}{10.9091}{mi}{112}\RG{5.488}{0.002}{10.9091}{rm}{40}\RG{9.731}{0.002}{10.9091}{mi}{65}\RG{17.913}{0.002}{10.9091}{rm}{41}\RG{25.168}{0.002}{10.9091}{rm}{61}\RG{36.677}{0.002}{10.9091}{rm}{48}%
}\hskip.1em\allowbreak{} ，所以

\[
\RB{98.342}{19.433}{15.217}{%
\RG{47.488}{13.633}{7.9701}{mi}{110}\RG{52.626}{13.633}{7.9701}{sy}{0}\RG{59.214}{13.633}{7.9701}{rm}{49}\RG{47.589}{10.360}{10.9092}{ex}{88}\RG{8.181}{4.500}{7.9701}{mi}{110}\RG{82.551}{4.500}{7.9701}{mi}{106}\RG{0.000}{-0.004}{10.9091}{mi}{65}\RG{13.818}{-0.004}{10.9091}{mi}{66}\RG{25.670}{-0.004}{10.9091}{rm}{61}\RG{37.185}{-0.004}{10.9091}{sy}{0}\RG{65.266}{-0.004}{10.9091}{mi}{99}\RG{74.369}{-0.004}{10.9091}{mi}{65}\RG{86.933}{-0.004}{10.9091}{mi}{66}\RG{95.747}{-0.004}{10.9091}{mi}{58}\RG{69.987}{-1.640}{7.9701}{mi}{106}\RG{48.115}{-13.091}{7.9701}{mi}{106}\RG{51.999}{-13.091}{7.9701}{rm}{61}\RG{58.583}{-13.091}{7.9701}{rm}{48}%
}
\]
因此\allowbreak{}\hskip.1em \RB{12.293}{8.677}{3.232}{%
\RG{0.000}{-0.005}{10.9091}{mi}{65}\RG{8.182}{-0.005}{10.9091}{rm}{40}%
}\hskip.1em\allowbreak{} Im\allowbreak{}\hskip.1em \RB{52.107}{9.730}{3.232}{%
\RG{0.000}{-0.005}{10.9091}{mi}{65}\RG{8.182}{3.954}{7.9701}{mi}{106}\RG{12.564}{-0.005}{10.9091}{mi}{66}\RG{21.386}{-0.005}{10.9091}{rm}{41}\RG{28.659}{-0.005}{10.9091}{sy}{18}\RG{40.174}{-0.005}{10.9091}{mi}{87}%
}\hskip.1em\allowbreak{} 对所有\allowbreak{}\hskip.1em \RB{64.614}{7.761}{2.741}{%
\RG{0.000}{-0.005}{10.9091}{rm}{48}\RG{8.484}{-0.005}{10.9091}{sy}{20}\RG{20.000}{-0.005}{10.9091}{mi}{106}\RG{28.147}{-0.005}{10.9091}{sy}{20}\RG{39.662}{-0.005}{10.9091}{mi}{110}\RG{48.634}{-0.005}{10.9091}{sy}{0}\RG{59.543}{-0.005}{10.9091}{rm}{49}%
}\hskip.1em\allowbreak{} 都成立，进而

\[
\RB{55.629}{8.684}{3.225}{%
\RG{0.000}{0.001}{10.9091}{mi}{65}\RG{8.182}{0.002}{10.9091}{rm}{40}\RG{12.424}{0.001}{10.9091}{mi}{87}\RG{24.242}{0.002}{10.9091}{rm}{41}\RG{31.515}{0.001}{10.9091}{sy}{18}\RG{43.030}{0.002}{10.9091}{mi}{87}\RG{53.034}{0.002}{10.9091}{mi}{58}%
}
\]
选取\allowbreak{}\hskip.1em \RB{11.933}{7.949}{0.742}{%
\RG{0.000}{-0.002}{10.9091}{mi}{87}%
}\hskip.1em\allowbreak{} 的一组基\allowbreak{}\hskip.1em \RB{43.733}{5.320}{2.607}{%
\RG{0.000}{-0.002}{10.9091}{mi}{118}\RG{5.288}{-1.639}{7.9701}{rm}{49}\RG{10.020}{-0.002}{10.9091}{mi}{59}\RG{14.853}{-0.002}{10.9091}{mi}{58}\RG{19.697}{-0.002}{10.9091}{mi}{58}\RG{24.540}{-0.002}{10.9091}{mi}{58}\RG{29.373}{-0.002}{10.9091}{mi}{59}\RG{34.217}{-0.002}{10.9091}{mi}{118}\RG{39.550}{-1.639}{7.9701}{mi}{114}%
}\hskip.1em\allowbreak{} ，扩充为\allowbreak{}\hskip.1em \RB{15.424}{7.971}{0.502}{%
\RG{0.000}{-0.002}{10.9091}{mi}{75}\RG{10.046}{3.956}{7.9701}{mi}{110}%
}\hskip.1em\allowbreak{} 的一组基\allowbreak{}\hskip.1em \RB{44.928}{5.320}{2.607}{%
\RG{0.000}{-0.002}{10.9091}{mi}{118}\RG{5.288}{-1.639}{7.9701}{rm}{49}\RG{10.020}{-0.002}{10.9091}{mi}{59}\RG{14.853}{-0.002}{10.9091}{mi}{58}\RG{19.697}{-0.002}{10.9091}{mi}{58}\RG{24.529}{-0.002}{10.9091}{mi}{58}\RG{29.373}{-0.002}{10.9091}{mi}{59}\RG{34.217}{-0.002}{10.9091}{mi}{118}\RG{39.550}{-1.639}{7.9701}{mi}{110}%
}\hskip.1em\allowbreak{} ，并令

\[
\RB{79.331}{8.686}{3.223}{%
\RG{0.000}{0.004}{10.9091}{mi}{80}\RG{11.549}{0.004}{10.9091}{rm}{61}\RG{23.047}{0.004}{10.9091}{rm}{40}\RG{27.307}{0.004}{10.9091}{mi}{118}\RG{37.327}{0.004}{10.9091}{mi}{59}\RG{42.160}{0.004}{10.9091}{mi}{58}\RG{47.004}{0.004}{10.9091}{mi}{58}\RG{51.837}{0.004}{10.9091}{mi}{58}\RG{56.680}{0.004}{10.9091}{mi}{59}\RG{61.524}{0.004}{10.9091}{mi}{118}\RG{72.494}{0.004}{10.9091}{rm}{41}\RG{76.737}{0.004}{10.9091}{mi}{58}\RG{32.594}{-1.633}{7.9701}{rm}{49}\RG{66.858}{-1.633}{7.9701}{mi}{110}%
}
\]
因为\allowbreak{}\hskip.1em \RB{11.933}{7.951}{0.740}{%
\RG{0.000}{0.000}{10.9091}{mi}{87}%
}\hskip.1em\allowbreak{} 是\allowbreak{}\hskip.1em \RB{8.366}{8.311}{0.500}{%
\RG{0.000}{0.000}{10.9091}{mi}{65}%
}\hskip.1em\allowbreak{} –不变子空间，故\allowbreak{}\hskip.1em \RB{42.117}{8.311}{2.217}{%
\RG{0.000}{0.000}{10.9091}{mi}{65}\RG{8.182}{0.000}{10.9091}{mi}{118}\RG{13.470}{-1.637}{7.9701}{mi}{105}\RG{19.881}{0.000}{10.9091}{sy}{50}\RG{30.184}{0.000}{10.9091}{mi}{87}%
}\hskip.1em\allowbreak{} （\allowbreak{}\hskip.1em \RB{43.561}{7.766}{1.995}{%
\RG{0.000}{0.000}{10.9091}{rm}{49}\RG{8.485}{0.000}{10.9091}{sy}{20}\RG{20.000}{0.000}{10.9091}{mi}{105}\RG{26.789}{0.000}{10.9091}{sy}{20}\RG{38.304}{0.000}{10.9091}{mi}{114}%
}\hskip.1em\allowbreak{} ）；又因为Im\allowbreak{}\hskip.1em \RB{35.300}{7.951}{1.995}{%
\RG{0.000}{0.000}{10.9091}{mi}{66}\RG{11.852}{0.000}{10.9091}{sy}{18}\RG{23.367}{0.000}{10.9091}{mi}{87}%
}\hskip.1em\allowbreak{} ，\allowbreak{}\hskip.1em \RB{8.747}{7.951}{0.500}{%
\RG{0.000}{0.000}{10.9091}{mi}{66}%
}\hskip.1em\allowbreak{} 的每一列均在\allowbreak{}\hskip.1em \RB{11.933}{7.951}{0.740}{%
\RG{0.000}{0.000}{10.9091}{mi}{87}%
}\hskip.1em\allowbreak{} 内。于是\allowbreak{}\hskip.1em \RB{36.746}{9.754}{0.505}{%
\RG{0.000}{-0.005}{10.9091}{mi}{80}\RG{8.519}{3.954}{7.9701}{sy}{0}\RG{15.106}{3.954}{7.9701}{rm}{49}\RG{19.838}{-0.005}{10.9091}{mi}{65}\RG{28.020}{-0.005}{10.9091}{mi}{80}%
}\hskip.1em\allowbreak{} 的左下角\allowbreak{}\hskip.1em \RB{52.014}{8.677}{3.232}{%
\RG{0.000}{-0.005}{10.9091}{rm}{40}\RG{4.242}{-0.005}{10.9091}{mi}{110}\RG{13.173}{-0.005}{10.9091}{sy}{0}\RG{24.040}{-0.005}{10.9091}{mi}{114}\RG{29.265}{-0.005}{10.9091}{rm}{41}\RG{35.889}{-0.005}{10.9091}{sy}{2}\RG{46.757}{-0.005}{10.9091}{mi}{114}%
}\hskip.1em\allowbreak{} 分块为零，\allowbreak{}\hskip.1em \RB{28.586}{9.754}{0.505}{%
\RG{0.000}{-0.005}{10.9091}{mi}{80}\RG{8.519}{3.954}{7.9701}{sy}{0}\RG{15.106}{3.954}{7.9701}{rm}{49}\RG{19.838}{-0.005}{10.9091}{mi}{66}%
}\hskip.1em\allowbreak{} 的下方\allowbreak{}\hskip.1em \RB{25.055}{5.317}{0.625}{%
\RG{0.000}{-0.005}{10.9091}{mi}{110}\RG{8.930}{-0.005}{10.9091}{sy}{0}\RG{19.798}{-0.005}{10.9091}{mi}{114}%
}\hskip.1em\allowbreak{} 行为零，即得所需形式。\allowbreak{}\hskip.1em \RB{25.272}{7.761}{0.745}{%
\RG{0.000}{-0.005}{10.9091}{mi}{114}\RG{8.255}{-0.005}{10.9091}{rm}{61}\RG{19.753}{-0.005}{10.9091}{rm}{48}%
}\hskip.1em\allowbreak{}或\allowbreak{}\hskip.1em \RB{26.499}{5.322}{0.620}{%
\RG{0.000}{0.000}{10.9091}{mi}{114}\RG{8.255}{0.000}{10.9091}{rm}{61}\RG{19.770}{0.000}{10.9091}{mi}{110}%
}\hskip.1em\allowbreak{} 时，空分块按通常约定处理。小结：由\allowbreak{}\hskip.1em \RB{88.303}{9.756}{2.608}{%
\RG{0.000}{-0.003}{10.9091}{mi}{66}\RG{8.815}{-0.003}{10.9091}{mi}{59}\RG{13.647}{-0.003}{10.9091}{mi}{65}\RG{21.829}{-0.003}{10.9091}{mi}{66}\RG{30.644}{-0.003}{10.9091}{mi}{59}\RG{35.488}{-0.003}{10.9091}{mi}{58}\RG{40.331}{-0.003}{10.9091}{mi}{58}\RG{45.164}{-0.003}{10.9091}{mi}{58}\RG{50.008}{-0.003}{10.9091}{mi}{59}\RG{54.851}{-0.003}{10.9091}{mi}{65}\RG{63.099}{3.956}{7.9701}{mi}{110}\RG{68.237}{3.956}{7.9701}{sy}{0}\RG{74.824}{3.956}{7.9701}{rm}{49}\RG{79.556}{-0.003}{10.9091}{mi}{66}%
}\hskip.1em\allowbreak{} 的列向量张成的空间在\allowbreak{}\hskip.1em \RB{8.366}{8.308}{0.503}{%
\RG{0.000}{-0.003}{10.9091}{mi}{65}%
}\hskip.1em\allowbreak{} 下不变；选择适应该不变子空间的基即可得到分块上三角形式。

\[
\RB{291.901}{7.727}{2.293}{%
\RT{0.000}{0.000}{7.9701}{题}\RT{8.966}{0.000}{7.9701}{目}\RT{17.933}{0.000}{7.9701}{来}\RT{26.899}{0.000}{7.9701}{源}\RT{35.866}{0.000}{7.9701}{：}\RT{44.832}{0.000}{7.9701}{谢}\RT{53.799}{0.000}{7.9701}{启}\RT{62.765}{0.000}{7.9701}{鸿}\RT{71.732}{0.000}{7.9701}{《}\RT{80.698}{0.000}{7.9701}{复}\RT{89.665}{0.000}{7.9701}{旦}\RT{98.631}{0.000}{7.9701}{大}\RT{107.598}{0.000}{7.9701}{学}\RT{116.564}{0.000}{7.9701}{高}\RT{125.531}{0.000}{7.9701}{等}\RT{134.497}{0.000}{7.9701}{代}\RT{143.464}{0.000}{7.9701}{数}\RT{152.430}{0.000}{7.9701}{习}\RT{161.397}{0.000}{7.9701}{题}\RT{170.363}{0.000}{7.9701}{集}\RT{179.329}{0.000}{7.9701}{（}\RT{188.296}{0.000}{7.9701}{2}\RT{192.779}{0.000}{7.9701}{0}\RT{197.262}{0.000}{7.9701}{2}\RT{201.745}{0.000}{7.9701}{3}\RT{209.215}{0.000}{7.9701}{版}\RT{218.181}{0.000}{7.9701}{）}\RT{222.664}{0.000}{7.9701}{》}\RT{227.148}{0.000}{7.9701}{，}\RT{236.114}{0.000}{7.9701}{第}\RT{248.066}{0.000}{7.9701}{7}\RT{252.549}{0.000}{7.9701}{4}\RT{257.033}{0.000}{7.9701}{–}\RT{261.516}{0.000}{7.9701}{7}\RT{265.999}{0.000}{7.9701}{5}\RT{273.468}{0.000}{7.9701}{页}\RT{282.435}{0.000}{7.9701}{。}%
}
\]
% Original page 7.
\section*{2019S07}
设\allowbreak{}\hskip.1em \RB{35.464}{8.313}{2.603}{%
\RG{0.000}{0.002}{10.9091}{mi}{65}\RG{8.182}{0.002}{10.9091}{mi}{59}\RG{13.015}{0.002}{10.9091}{mi}{66}\RG{21.829}{0.002}{10.9091}{mi}{59}\RG{26.673}{0.002}{10.9091}{mi}{67}%
}\hskip.1em\allowbreak{} 为\allowbreak{}\hskip.1em \RB{6.729}{5.324}{0.618}{%
\RG{0.000}{0.002}{10.9091}{mi}{110}%
}\hskip.1em\allowbreak{} 阶复矩阵，满足\allowbreak{}\hskip.1em \RB{70.641}{8.313}{0.738}{%
\RG{0.000}{0.002}{10.9091}{mi}{67}\RG{11.608}{0.002}{10.9091}{rm}{61}\RG{23.123}{0.002}{10.9091}{mi}{65}\RG{31.305}{0.002}{10.9091}{mi}{66}\RG{42.551}{0.002}{10.9091}{sy}{0}\RG{53.460}{0.002}{10.9091}{mi}{66}\RG{62.275}{0.002}{10.9091}{mi}{65}%
}\hskip.1em\allowbreak{} 、\allowbreak{}\hskip.1em \RB{48.233}{8.313}{0.738}{%
\RG{0.000}{0.002}{10.9091}{mi}{65}\RG{8.182}{0.002}{10.9091}{mi}{67}\RG{19.789}{0.002}{10.9091}{rm}{61}\RG{31.304}{0.002}{10.9091}{mi}{67}\RG{39.868}{0.002}{10.9091}{mi}{65}%
}\hskip.1em\allowbreak{} 、\allowbreak{}\hskip.1em \RB{49.255}{8.193}{0.738}{%
\RG{0.000}{0.002}{10.9091}{mi}{66}\RG{8.815}{0.002}{10.9091}{mi}{67}\RG{20.429}{0.002}{10.9091}{rm}{61}\RG{31.944}{0.002}{10.9091}{mi}{67}\RG{40.508}{0.002}{10.9091}{mi}{66}%
}\hskip.1em\allowbreak{} 。（1）利用Jordan标准型理论证明\allowbreak{}\hskip.1em \RB{8.791}{8.190}{0.741}{%
\RG{0.000}{-0.001}{10.9091}{mi}{67}%
}\hskip.1em\allowbreak{} 的特征值全为零。（2）设\allowbreak{}\hskip.1em \RB{37.683}{5.319}{2.608}{%
\RG{0.000}{-0.003}{10.9091}{mi}{109}\RG{9.578}{-1.692}{7.9701}{mi}{65}\RG{16.419}{-0.003}{10.9091}{mi}{59}\RG{21.252}{-0.003}{10.9091}{mi}{109}\RG{30.846}{-1.692}{7.9701}{mi}{66}%
}\hskip.1em\allowbreak{} 分别为\allowbreak{}\hskip.1em \RB{21.762}{8.308}{2.608}{%
\RG{0.000}{-0.003}{10.9091}{mi}{65}\RG{8.182}{-0.003}{10.9091}{mi}{59}\RG{13.015}{-0.003}{10.9091}{mi}{66}%
}\hskip.1em\allowbreak{} 的极小多项式，令\allowbreak{}\hskip.1em \RB{17.419}{8.068}{0.623}{%
\RG{0.000}{-0.003}{10.9091}{mi}{107}\RG{9.053}{-0.003}{10.9091}{rm}{61}%
}\hskip.1em\allowbreak{} min\allowbreak{}\hskip.1em \RB{5.158}{8.679}{3.230}{%
\RG{0.000}{-0.003}{10.9091}{sy}{102}%
}\hskip.1em\allowbreak{} deg\allowbreak{}\hskip.1em \RB{19.135}{5.319}{2.608}{%
\RG{0.000}{-0.003}{10.9091}{mi}{109}\RG{9.579}{-1.692}{7.9701}{mi}{65}\RG{16.420}{-0.003}{10.9091}{mi}{59}%
}\hskip.1em\allowbreak{} deg\allowbreak{}\hskip.1em \RB{49.711}{8.679}{3.230}{%
\RG{0.000}{-0.003}{10.9091}{mi}{109}\RG{9.578}{-1.692}{7.9701}{mi}{66}\RG{16.857}{-0.003}{10.9091}{mi}{59}\RG{21.690}{-0.003}{10.9091}{mi}{110}\RG{29.433}{-0.003}{10.9091}{sy}{0}\RG{39.098}{-0.003}{10.9091}{rm}{49}\RG{44.552}{-0.003}{10.9091}{sy}{103}%
}\hskip.1em\allowbreak{} 。证明\allowbreak{}\hskip.1em \RB{33.744}{9.987}{0.743}{%
\RG{0.000}{-0.003}{10.9091}{mi}{67}\RG{8.577}{3.956}{7.9701}{mi}{107}\RG{16.727}{-0.003}{10.9091}{rm}{61}\RG{28.225}{-0.003}{10.9091}{rm}{48}%
}\hskip.1em\allowbreak{} 。

证明：（1）由\allowbreak{}\hskip.1em \RB{48.234}{8.316}{0.735}{%
\RG{0.000}{0.005}{10.9091}{mi}{65}\RG{8.182}{0.005}{10.9091}{mi}{67}\RG{19.790}{0.005}{10.9091}{rm}{61}\RG{31.305}{0.005}{10.9091}{mi}{67}\RG{39.869}{0.005}{10.9091}{mi}{65}%
}\hskip.1em\allowbreak{} 、\allowbreak{}\hskip.1em \RB{49.255}{8.196}{0.735}{%
\RG{0.000}{0.005}{10.9091}{mi}{66}\RG{8.815}{0.005}{10.9091}{mi}{67}\RG{20.429}{0.005}{10.9091}{rm}{61}\RG{31.944}{0.005}{10.9091}{mi}{67}\RG{40.508}{0.005}{10.9091}{mi}{66}%
}\hskip.1em\allowbreak{} 及迹的轮换性质，对任意\allowbreak{}\hskip.1em \RB{24.733}{7.771}{2.731}{%
\RG{0.000}{0.005}{10.9091}{mi}{106}\RG{8.147}{0.005}{10.9091}{sy}{21}\RG{19.662}{0.005}{10.9091}{rm}{49}%
}\hskip.1em\allowbreak{} ，

\[
\RB{196.017}{28.465}{3.230}{%
\RG{53.030}{26.997}{10.9092}{ex}{0}\RG{137.637}{26.997}{10.9092}{ex}{1}\RG{21.340}{22.665}{7.9701}{mi}{106}\RG{122.434}{22.665}{7.9701}{mi}{106}\RG{126.318}{22.665}{7.9701}{sy}{0}\RG{132.905}{22.665}{7.9701}{rm}{49}\RT{0.000}{18.161}{10.9091}{t}\RT{4.244}{18.161}{10.9091}{r}\RG{8.520}{18.161}{10.9091}{rm}{40}\RG{12.763}{18.161}{10.9091}{mi}{67}\RG{25.722}{18.161}{10.9091}{rm}{41}\RG{32.977}{18.161}{10.9091}{rm}{61}\RT{44.510}{18.161}{10.9091}{t}\RT{48.754}{18.161}{10.9091}{r}\RG{58.030}{18.161}{10.9091}{rm}{40}\RG{62.274}{18.161}{10.9091}{mi}{65}\RG{70.455}{18.161}{10.9091}{mi}{66}\RG{81.702}{18.161}{10.9091}{sy}{0}\RG{92.611}{18.161}{10.9091}{mi}{66}\RG{101.425}{18.161}{10.9091}{mi}{65}\RG{109.614}{18.161}{10.9091}{rm}{41}\RG{113.857}{18.161}{10.9091}{mi}{67}\RG{82.855}{4.501}{7.9701}{mi}{106}\RG{86.739}{4.501}{7.9701}{sy}{0}\RG{93.325}{4.501}{7.9701}{rm}{49}\RG{145.155}{4.501}{7.9701}{mi}{106}\RG{149.039}{4.501}{7.9701}{sy}{0}\RG{155.625}{4.501}{7.9701}{rm}{49}\RG{32.995}{-0.003}{10.9091}{rm}{61}\RT{44.510}{-0.003}{10.9091}{t}\RT{48.754}{-0.003}{10.9091}{r}\RG{53.030}{-0.003}{10.9091}{rm}{40}\RG{57.273}{-0.003}{10.9091}{mi}{65}\RG{65.455}{-0.003}{10.9091}{mi}{66}\RG{74.270}{-0.003}{10.9091}{mi}{67}\RG{98.057}{-0.003}{10.9091}{rm}{41}\RG{104.724}{-0.003}{10.9091}{sy}{0}\RT{115.633}{-0.003}{10.9091}{t}\RT{119.877}{-0.003}{10.9091}{r}\RG{124.153}{-0.003}{10.9091}{rm}{40}\RG{128.396}{-0.003}{10.9091}{mi}{65}\RG{136.578}{-0.003}{10.9091}{mi}{67}\RG{160.358}{-0.003}{10.9091}{mi}{66}\RG{169.180}{-0.003}{10.9091}{rm}{41}\RG{176.435}{-0.003}{10.9091}{rm}{61}\RG{187.944}{-0.003}{10.9091}{rm}{48}\RG{193.423}{-0.003}{10.9091}{mi}{58}%
}
\]
设\allowbreak{}\hskip.1em \RB{47.080}{8.074}{2.602}{%
\RG{0.000}{0.003}{10.9091}{mi}{21}\RG{6.363}{-1.633}{7.9701}{rm}{49}\RG{11.096}{0.003}{10.9091}{mi}{59}\RG{15.929}{0.003}{10.9091}{mi}{58}\RG{20.773}{0.003}{10.9091}{mi}{58}\RG{25.605}{0.003}{10.9091}{mi}{58}\RG{30.449}{0.003}{10.9091}{mi}{59}\RG{35.293}{0.003}{10.9091}{mi}{21}\RG{41.702}{-1.633}{7.9701}{mi}{110}%
}\hskip.1em\allowbreak{} 为\allowbreak{}\hskip.1em \RB{8.791}{8.194}{0.737}{%
\RG{0.000}{0.003}{10.9091}{mi}{67}%
}\hskip.1em\allowbreak{} 的特征值，计入代数重数。根据Jordan标准型，

\[
\RB{171.941}{17.649}{13.669}{%
\RG{5.310}{13.634}{7.9701}{mi}{110}\RG{0.000}{10.362}{10.9092}{ex}{88}\RG{23.940}{5.305}{7.9701}{mi}{106}\RG{64.208}{4.502}{7.9701}{mi}{106}\RG{17.576}{-0.002}{10.9091}{mi}{21}\RG{31.352}{-0.002}{10.9091}{rm}{61}\RT{42.867}{-0.002}{10.9091}{t}\RT{47.111}{-0.002}{10.9091}{r}\RG{51.387}{-0.002}{10.9091}{rm}{40}\RG{55.630}{-0.002}{10.9091}{mi}{67}\RG{68.590}{-0.002}{10.9091}{rm}{41}\RG{75.844}{-0.002}{10.9091}{rm}{61}\RG{87.353}{-0.002}{10.9091}{rm}{48}\RG{114.626}{-0.002}{10.9091}{rm}{40}\RG{118.859}{-0.002}{10.9091}{rm}{49}\RG{127.378}{-0.002}{10.9091}{sy}{20}\RG{138.893}{-0.002}{10.9091}{mi}{106}\RG{147.040}{-0.002}{10.9091}{sy}{20}\RG{158.555}{-0.002}{10.9091}{mi}{110}\RG{165.103}{-0.002}{10.9091}{rm}{41}\RG{169.346}{-0.002}{10.9091}{mi}{58}\RG{23.940}{-3.259}{7.9701}{mi}{105}\RG{1.027}{-13.089}{7.9701}{mi}{105}\RG{3.911}{-13.089}{7.9701}{rm}{61}\RG{10.494}{-13.089}{7.9701}{rm}{49}%
}
\]
由Newton恒等式，\allowbreak{}\hskip.1em \RB{47.080}{8.075}{2.601}{%
\RG{0.000}{0.004}{10.9091}{mi}{21}\RG{6.364}{-1.633}{7.9701}{rm}{49}\RG{11.096}{0.004}{10.9091}{mi}{59}\RG{15.929}{0.004}{10.9091}{mi}{58}\RG{20.773}{0.004}{10.9091}{mi}{58}\RG{25.616}{0.004}{10.9091}{mi}{58}\RG{30.449}{0.004}{10.9091}{mi}{59}\RG{35.293}{0.004}{10.9091}{mi}{21}\RG{41.702}{-1.633}{7.9701}{mi}{110}%
}\hskip.1em\allowbreak{} 的全部正次数初等对称多项式均为零，因而det\allowbreak{}\hskip.1em \RB{72.119}{8.686}{3.223}{%
\RG{0.000}{0.004}{10.9091}{rm}{40}\RG{4.242}{0.004}{10.9091}{mi}{120}\RG{10.471}{0.004}{10.9091}{mi}{73}\RG{15.273}{-1.633}{7.9701}{mi}{110}\RG{22.782}{0.004}{10.9091}{sy}{0}\RG{33.141}{0.004}{10.9091}{mi}{67}\RG{41.718}{0.004}{10.9091}{rm}{41}\RG{48.973}{0.004}{10.9091}{rm}{61}\RG{60.507}{0.004}{10.9091}{mi}{120}\RG{66.742}{3.963}{7.9701}{mi}{110}%
}\hskip.1em\allowbreak{} 。故\allowbreak{}\hskip.1em \RB{8.791}{8.190}{0.741}{%
\RG{0.000}{-0.001}{10.9091}{mi}{67}%
}\hskip.1em\allowbreak{} 的特征值全为零，且\allowbreak{}\hskip.1em \RB{8.791}{8.190}{0.741}{%
\RG{0.000}{-0.001}{10.9091}{mi}{67}%
}\hskip.1em\allowbreak{} 幂零。（2）设\allowbreak{}\hskip.1em \RB{23.915}{8.067}{2.192}{%
\RG{0.000}{-0.004}{10.9091}{mi}{100}\RG{5.678}{-1.692}{7.9701}{mi}{65}\RG{15.549}{-0.004}{10.9091}{rm}{61}%
}\hskip.1em\allowbreak{} deg\allowbreak{}\hskip.1em \RB{16.167}{5.318}{2.192}{%
\RG{0.000}{-0.004}{10.9091}{mi}{109}\RG{9.578}{-1.692}{7.9701}{mi}{65}%
}\hskip.1em\allowbreak{} ，定义\allowbreak{}\hskip.1em \RB{92.767}{8.678}{3.231}{%
\RG{0.000}{-0.004}{10.9091}{mi}{14}\RG{4.849}{-1.692}{7.9701}{mi}{66}\RG{12.127}{-0.004}{10.9091}{rm}{40}\RG{16.369}{-0.004}{10.9091}{mi}{88}\RG{26.263}{-0.004}{10.9091}{rm}{41}\RG{33.518}{-0.004}{10.9091}{rm}{61}\RG{45.051}{-0.004}{10.9091}{mi}{88}\RG{54.935}{-0.004}{10.9091}{mi}{66}\RG{64.726}{-0.004}{10.9091}{sy}{0}\RG{74.169}{-0.004}{10.9091}{mi}{66}\RG{82.983}{-0.004}{10.9091}{mi}{88}%
}\hskip.1em\allowbreak{} 。由题设\allowbreak{}\hskip.1em \RB{52.130}{8.678}{3.231}{%
\RG{0.000}{-0.004}{10.9091}{mi}{14}\RG{4.849}{-1.692}{7.9701}{mi}{66}\RG{12.127}{-0.004}{10.9091}{rm}{40}\RG{16.370}{-0.004}{10.9091}{mi}{65}\RG{24.551}{-0.004}{10.9091}{rm}{41}\RG{31.806}{-0.004}{10.9091}{rm}{61}\RG{43.339}{-0.004}{10.9091}{mi}{67}%
}\hskip.1em\allowbreak{} ，\allowbreak{}\hskip.1em \RB{49.229}{8.678}{3.231}{%
\RG{0.000}{-0.004}{10.9091}{mi}{14}\RG{4.849}{-1.692}{7.9701}{mi}{66}\RG{12.127}{-0.004}{10.9091}{rm}{40}\RG{16.370}{-0.004}{10.9091}{mi}{67}\RG{24.947}{-0.004}{10.9091}{rm}{41}\RG{32.202}{-0.004}{10.9091}{rm}{61}\RG{43.711}{-0.004}{10.9091}{rm}{48}%
}\hskip.1em\allowbreak{} ，且\allowbreak{}\hskip.1em \RB{48.233}{8.307}{0.744}{%
\RG{0.000}{-0.004}{10.9091}{mi}{65}\RG{8.182}{-0.004}{10.9091}{mi}{67}\RG{19.789}{-0.004}{10.9091}{rm}{61}\RG{31.304}{-0.004}{10.9091}{mi}{67}\RG{39.868}{-0.004}{10.9091}{mi}{65}%
}\hskip.1em\allowbreak{} 。用Leibniz法则归纳可得

\[
\RB{160.281}{11.080}{3.938}{%
\RG{5.262}{5.304}{7.9701}{mi}{106}\RG{64.856}{4.501}{7.9701}{rm}{40}\RG{68.149}{4.501}{7.9701}{mi}{106}\RG{72.033}{4.501}{7.9701}{rm}{41}\RG{101.068}{4.501}{7.9701}{mi}{106}\RG{0.000}{-0.004}{10.9091}{mi}{14}\RG{12.127}{-0.004}{10.9091}{rm}{40}\RG{16.370}{-0.004}{10.9091}{mi}{102}\RG{22.885}{-0.004}{10.9091}{rm}{40}\RG{27.127}{-0.004}{10.9091}{mi}{65}\RG{35.309}{-0.004}{10.9091}{rm}{41}\RG{39.542}{-0.004}{10.9091}{rm}{41}\RG{46.797}{-0.004}{10.9091}{rm}{61}\RG{58.339}{-0.004}{10.9091}{mi}{102}\RG{75.824}{-0.004}{10.9091}{rm}{40}\RG{80.067}{-0.004}{10.9091}{mi}{65}\RG{88.249}{-0.004}{10.9091}{rm}{41}\RG{92.491}{-0.004}{10.9091}{mi}{67}\RG{127.269}{-0.004}{10.9091}{rm}{40}\RG{131.511}{-0.004}{10.9091}{mi}{106}\RG{139.658}{-0.004}{10.9091}{sy}{21}\RG{151.173}{-0.004}{10.9091}{rm}{48}\RG{156.628}{-0.004}{10.9091}{rm}{41}\RG{4.849}{-3.438}{7.9701}{mi}{66}%
}
\]
对任意多项式\allowbreak{}\hskip.1em \RB{6.522}{8.187}{2.740}{%
\RG{0.000}{-0.004}{10.9091}{mi}{102}%
}\hskip.1em\allowbreak{} 都成立。由\allowbreak{}\hskip.1em \RB{53.126}{8.678}{3.231}{%
\RG{0.000}{-0.004}{10.9091}{mi}{109}\RG{9.579}{-1.693}{7.9701}{mi}{65}\RG{16.420}{-0.004}{10.9091}{rm}{40}\RG{20.662}{-0.004}{10.9091}{mi}{65}\RG{28.844}{-0.004}{10.9091}{rm}{41}\RG{36.099}{-0.004}{10.9091}{rm}{61}\RG{47.608}{-0.004}{10.9091}{rm}{48}%
}\hskip.1em\allowbreak{} 且\allowbreak{}\hskip.1em \RB{16.167}{5.318}{2.193}{%
\RG{0.000}{-0.004}{10.9091}{mi}{109}\RG{9.578}{-1.693}{7.9701}{mi}{65}%
}\hskip.1em\allowbreak{} 首一，

\[
\RB{131.793}{11.130}{3.932}{%
\RG{25.261}{5.099}{7.9701}{mi}{100}\RG{118.201}{4.506}{7.9701}{mi}{100}\RG{29.619}{3.759}{5.9776}{mi}{65}\RG{122.558}{3.166}{5.9776}{mi}{65}\RG{0.000}{0.002}{10.9091}{rm}{48}\RG{8.476}{0.002}{10.9091}{rm}{61}\RG{20.000}{0.002}{10.9091}{mi}{14}\RG{36.139}{0.002}{10.9091}{rm}{40}\RG{40.381}{0.002}{10.9091}{mi}{109}\RG{56.801}{0.002}{10.9091}{rm}{40}\RG{61.044}{0.002}{10.9091}{mi}{65}\RG{69.226}{0.002}{10.9091}{rm}{41}\RG{73.458}{0.002}{10.9091}{rm}{41}\RG{80.713}{0.002}{10.9091}{rm}{61}\RG{92.256}{0.002}{10.9091}{mi}{100}\RG{104.775}{0.002}{10.9091}{rm}{33}\RG{109.624}{0.002}{10.9091}{mi}{67}\RG{129.078}{0.002}{10.9091}{mi}{59}\RG{49.960}{-1.687}{7.9701}{mi}{65}\RG{97.934}{-1.687}{7.9701}{mi}{65}\RG{24.849}{-3.432}{7.9701}{mi}{66}%
}
\]
因此\allowbreak{}\hskip.1em \RB{39.742}{9.988}{0.742}{%
\RG{0.000}{-0.002}{10.9091}{mi}{67}\RG{8.577}{3.957}{7.9701}{mi}{100}\RG{12.935}{2.617}{5.9776}{mi}{65}\RG{22.605}{-0.002}{10.9091}{rm}{61}\RG{34.223}{-0.002}{10.9091}{rm}{48}%
}\hskip.1em\allowbreak{} 。类似地，定义\allowbreak{}\hskip.1em \RB{94.610}{8.680}{3.229}{%
\RG{0.000}{-0.002}{10.9091}{mi}{14}\RG{4.849}{-1.691}{7.9701}{mi}{65}\RG{11.690}{-0.002}{10.9091}{rm}{40}\RG{15.932}{-0.002}{10.9091}{mi}{88}\RG{25.826}{-0.002}{10.9091}{rm}{41}\RG{33.201}{-0.002}{10.9091}{rm}{61}\RG{44.854}{-0.002}{10.9091}{mi}{65}\RG{53.036}{-0.002}{10.9091}{mi}{88}\RG{65.403}{-0.002}{10.9091}{sy}{0}\RG{76.361}{-0.002}{10.9091}{mi}{88}\RG{86.244}{-0.002}{10.9091}{mi}{65}%
}\hskip.1em\allowbreak{} ，利用\allowbreak{}\hskip.1em \RB{52.573}{8.680}{3.229}{%
\RG{0.000}{-0.002}{10.9091}{mi}{14}\RG{4.848}{-1.691}{7.9701}{mi}{65}\RG{11.689}{-0.002}{10.9091}{rm}{40}\RG{15.932}{-0.002}{10.9091}{mi}{66}\RG{24.754}{-0.002}{10.9091}{rm}{41}\RG{32.129}{-0.002}{10.9091}{rm}{61}\RG{43.782}{-0.002}{10.9091}{mi}{67}%
}\hskip.1em\allowbreak{} 、\allowbreak{}\hskip.1em \RB{49.032}{8.680}{3.229}{%
\RG{0.000}{-0.002}{10.9091}{mi}{14}\RG{4.849}{-1.691}{7.9701}{mi}{65}\RG{11.690}{-0.002}{10.9091}{rm}{40}\RG{15.932}{-0.002}{10.9091}{mi}{67}\RG{24.510}{-0.002}{10.9091}{rm}{41}\RG{31.885}{-0.002}{10.9091}{rm}{61}\RG{43.514}{-0.002}{10.9091}{rm}{48}%
}\hskip.1em\allowbreak{} 与\allowbreak{}\hskip.1em \RB{54.443}{8.680}{3.229}{%
\RG{0.000}{-0.002}{10.9091}{mi}{109}\RG{9.578}{-1.691}{7.9701}{mi}{66}\RG{16.857}{-0.002}{10.9091}{rm}{40}\RG{21.099}{-0.002}{10.9091}{mi}{66}\RG{29.921}{-0.002}{10.9091}{rm}{41}\RG{37.296}{-0.002}{10.9091}{rm}{61}\RG{48.925}{-0.002}{10.9091}{rm}{48}%
}\hskip.1em\allowbreak{} ，得到\allowbreak{}\hskip.1em \RB{39.803}{9.992}{0.737}{%
\RG{0.000}{0.003}{10.9091}{mi}{67}\RG{8.577}{3.961}{7.9701}{mi}{100}\RG{12.934}{2.622}{5.9776}{mi}{66}\RG{22.786}{0.003}{10.9091}{rm}{61}\RG{34.284}{0.003}{10.9091}{rm}{48}%
}\hskip.1em\allowbreak{} ，其中\allowbreak{}\hskip.1em \RB{24.353}{8.074}{2.186}{%
\RG{0.000}{0.003}{10.9091}{mi}{100}\RG{5.678}{-1.686}{7.9701}{mi}{66}\RG{15.987}{0.003}{10.9091}{rm}{61}%
}\hskip.1em\allowbreak{} deg\allowbreak{}\hskip.1em \RB{16.414}{5.325}{2.186}{%
\RG{0.000}{0.003}{10.9091}{mi}{109}\RG{9.578}{-1.686}{7.9701}{mi}{66}%
}\hskip.1em\allowbreak{} 。

再证\allowbreak{}\hskip.1em \RB{45.189}{9.759}{0.740}{%
\RG{0.000}{0.000}{10.9091}{mi}{67}\RG{8.578}{3.959}{7.9701}{mi}{110}\RG{13.716}{3.959}{7.9701}{sy}{0}\RG{20.302}{3.959}{7.9701}{rm}{49}\RG{28.118}{0.000}{10.9091}{rm}{61}\RG{39.671}{0.000}{10.9091}{rm}{48}%
}\hskip.1em\allowbreak{} （\allowbreak{}\hskip.1em \RB{26.598}{7.766}{1.995}{%
\RG{0.000}{0.000}{10.9091}{mi}{110}\RG{9.632}{0.000}{10.9091}{sy}{21}\RG{21.200}{0.000}{10.9091}{rm}{50}%
}\hskip.1em\allowbreak{} ）。若\allowbreak{}\hskip.1em \RB{45.189}{9.759}{2.856}{%
\RG{0.000}{0.000}{10.9091}{mi}{67}\RG{8.577}{3.959}{7.9701}{mi}{110}\RG{13.715}{3.959}{7.9701}{sy}{0}\RG{20.302}{3.959}{7.9701}{rm}{49}\RG{28.118}{0.000}{10.9091}{sy}{54}\RG{28.118}{0.000}{10.9091}{rm}{61}\RG{39.671}{0.000}{10.9091}{rm}{48}%
}\hskip.1em\allowbreak{} ，由（1）可知\allowbreak{}\hskip.1em \RB{8.791}{8.191}{0.740}{%
\RG{0.000}{0.000}{10.9091}{mi}{67}%
}\hskip.1em\allowbreak{} 的Jordan标准型必为单个块\allowbreak{}\hskip.1em \RB{25.026}{8.682}{3.227}{%
\RG{0.000}{0.000}{10.9091}{mi}{74}\RG{6.049}{-1.636}{7.9701}{mi}{110}\RG{11.686}{0.000}{10.9091}{rm}{40}\RG{15.919}{0.000}{10.9091}{rm}{48}\RG{21.373}{0.000}{10.9091}{rm}{41}%
}\hskip.1em\allowbreak{} 。故存在向量\allowbreak{}\hskip.1em \RB{5.606}{5.317}{0.625}{%
\RG{0.000}{-0.005}{10.9091}{mi}{118}%
}\hskip.1em\allowbreak{} ，使得\allowbreak{}\hskip.1em \RB{79.667}{9.754}{2.610}{%
\RG{0.000}{-0.005}{10.9091}{mi}{118}\RG{5.662}{-0.005}{10.9091}{mi}{59}\RG{10.506}{-0.005}{10.9091}{mi}{67}\RG{19.080}{-0.005}{10.9091}{mi}{118}\RG{24.742}{-0.005}{10.9091}{mi}{59}\RG{29.586}{-0.005}{10.9091}{mi}{58}\RG{34.429}{-0.005}{10.9091}{mi}{58}\RG{39.262}{-0.005}{10.9091}{mi}{58}\RG{44.106}{-0.005}{10.9091}{mi}{59}\RG{48.950}{-0.005}{10.9091}{mi}{67}\RG{57.604}{3.954}{7.9701}{mi}{110}\RG{62.743}{3.954}{7.9701}{sy}{0}\RG{69.330}{3.954}{7.9701}{rm}{49}\RG{74.062}{-0.005}{10.9091}{mi}{118}%
}\hskip.1em\allowbreak{} 成为一组基。对任一与\allowbreak{}\hskip.1em \RB{8.791}{8.186}{0.745}{%
\RG{0.000}{-0.005}{10.9091}{mi}{67}%
}\hskip.1em\allowbreak{} 交换的矩阵\allowbreak{}\hskip.1em \RB{8.180}{7.881}{0.505}{%
\RG{0.000}{-0.005}{10.9091}{mi}{84}%
}\hskip.1em\allowbreak{} ，可将\allowbreak{}\hskip.1em \RB{13.482}{7.881}{0.625}{%
\RG{0.000}{-0.005}{10.9091}{mi}{84}\RG{7.876}{-0.005}{10.9091}{mi}{118}%
}\hskip.1em\allowbreak{} 唯一写为\allowbreak{}\hskip.1em \RB{28.157}{8.682}{3.227}{%
\RG{0.000}{0.000}{10.9091}{mi}{112}\RG{5.488}{0.000}{10.9091}{rm}{40}\RG{9.731}{0.000}{10.9091}{mi}{67}\RG{18.308}{0.000}{10.9091}{rm}{41}\RG{22.551}{0.000}{10.9091}{mi}{118}%
}\hskip.1em\allowbreak{} ，其中deg\allowbreak{}\hskip.1em \RB{26.748}{6.391}{2.616}{%
\RG{0.000}{0.000}{10.9091}{mi}{112}\RG{8.509}{0.000}{10.9091}{mi}{60}\RG{20.018}{0.000}{10.9091}{mi}{110}%
}\hskip.1em\allowbreak{} ，于是

\[
\RB{202.567}{10.277}{3.231}{%
\RG{16.468}{4.501}{7.9701}{mi}{106}\RG{49.652}{4.501}{7.9701}{mi}{106}\RG{113.278}{4.501}{7.9701}{mi}{106}\RG{0.000}{-0.004}{10.9091}{mi}{84}\RG{7.876}{-0.004}{10.9091}{mi}{67}\RG{20.850}{-0.004}{10.9091}{mi}{118}\RG{29.559}{-0.004}{10.9091}{rm}{61}\RG{41.074}{-0.004}{10.9091}{mi}{67}\RG{54.034}{-0.004}{10.9091}{mi}{84}\RG{61.911}{-0.004}{10.9091}{mi}{118}\RG{70.635}{-0.004}{10.9091}{rm}{61}\RG{82.150}{-0.004}{10.9091}{mi}{112}\RG{87.638}{-0.004}{10.9091}{rm}{40}\RG{91.881}{-0.004}{10.9091}{mi}{67}\RG{100.458}{-0.004}{10.9091}{rm}{41}\RG{104.700}{-0.004}{10.9091}{mi}{67}\RG{117.660}{-0.004}{10.9091}{mi}{118}\RG{145.157}{-0.004}{10.9091}{rm}{40}\RG{149.390}{-0.004}{10.9091}{rm}{48}\RG{157.885}{-0.004}{10.9091}{sy}{20}\RG{169.400}{-0.004}{10.9091}{mi}{106}\RG{177.528}{-0.004}{10.9091}{mi}{60}\RG{189.037}{-0.004}{10.9091}{mi}{110}\RG{195.611}{-0.004}{10.9091}{rm}{41}\RG{199.853}{-0.004}{10.9091}{mi}{59}%
}
\]
故\allowbreak{}\hskip.1em \RB{44.396}{8.684}{3.225}{%
\RG{0.000}{0.002}{10.9091}{mi}{84}\RG{10.920}{0.002}{10.9091}{rm}{61}\RG{22.435}{0.002}{10.9091}{mi}{112}\RG{27.924}{0.002}{10.9091}{rm}{40}\RG{32.166}{0.002}{10.9091}{mi}{67}\RG{40.743}{0.002}{10.9091}{rm}{41}%
}\hskip.1em\allowbreak{} 。特别地，\allowbreak{}\hskip.1em \RB{21.762}{8.313}{2.603}{%
\RG{0.000}{0.002}{10.9091}{mi}{65}\RG{8.182}{0.002}{10.9091}{mi}{59}\RG{13.015}{0.002}{10.9091}{mi}{66}%
}\hskip.1em\allowbreak{} 都是\allowbreak{}\hskip.1em \RB{8.791}{8.193}{0.738}{%
\RG{0.000}{0.002}{10.9091}{mi}{67}%
}\hskip.1em\allowbreak{} 的多项式，从而\allowbreak{}\hskip.1em \RB{48.729}{8.313}{0.500}{%
\RG{0.000}{0.002}{10.9091}{mi}{65}\RG{8.182}{0.002}{10.9091}{mi}{66}\RG{20.034}{0.002}{10.9091}{rm}{61}\RG{31.549}{0.002}{10.9091}{mi}{66}\RG{40.364}{0.002}{10.9091}{mi}{65}%
}\hskip.1em\allowbreak{} ，即\allowbreak{}\hskip.1em \RB{28.624}{8.193}{0.738}{%
\RG{0.000}{0.002}{10.9091}{mi}{67}\RG{11.607}{0.002}{10.9091}{rm}{61}\RG{23.105}{0.002}{10.9091}{rm}{48}%
}\hskip.1em\allowbreak{} ，矛盾。因此\allowbreak{}\hskip.1em \RB{45.081}{9.761}{0.738}{%
\RG{0.000}{0.002}{10.9091}{mi}{67}\RG{8.577}{3.961}{7.9701}{mi}{110}\RG{13.715}{3.961}{7.9701}{sy}{0}\RG{20.302}{3.961}{7.9701}{rm}{49}\RG{28.064}{0.002}{10.9091}{rm}{61}\RG{39.563}{0.002}{10.9091}{rm}{48}%
}\hskip.1em\allowbreak{} ，综合三个幂次界即得

\[
\RB{184.999}{6.928}{7.728}{%
\RT{8.577}{0.004}{7.9701}{m}\RT{15.216}{0.004}{7.9701}{i}\RT{17.432}{0.004}{7.9701}{n}\RG{21.864}{0.004}{7.9701}{sy}{102}\RT{26.098}{0.004}{7.9701}{d}\RT{30.530}{0.004}{7.9701}{e}\RT{34.068}{0.004}{7.9701}{g}\RG{39.464}{0.004}{7.9701}{mi}{109}\RG{52.977}{0.004}{7.9701}{mi}{59}\RT{55.329}{0.004}{7.9701}{d}\RT{59.761}{0.004}{7.9701}{e}\RT{63.300}{0.004}{7.9701}{g}\RG{68.697}{0.004}{7.9701}{mi}{109}\RG{82.510}{0.004}{7.9701}{mi}{59}\RG{84.861}{0.004}{7.9701}{mi}{110}\RG{90.000}{0.004}{7.9701}{sy}{0}\RG{96.587}{0.004}{7.9701}{rm}{49}\RG{100.821}{0.004}{7.9701}{sy}{103}\RG{46.955}{-1.336}{5.9776}{mi}{65}\RG{76.187}{-1.336}{5.9776}{mi}{66}\RG{0.000}{-4.501}{10.9091}{mi}{67}\RG{108.584}{-4.501}{10.9091}{rm}{61}\RG{120.082}{-4.501}{10.9091}{rm}{48}\RG{147.355}{-4.501}{10.9091}{rm}{40}\RG{151.614}{-4.501}{10.9091}{mi}{110}\RG{161.192}{-4.501}{10.9091}{sy}{21}\RG{172.707}{-4.501}{10.9091}{rm}{50}\RG{178.162}{-4.501}{10.9091}{rm}{41}\RG{182.405}{-4.501}{10.9091}{mi}{58}%
}
\]
边界说明：原题第（2）问未写明\allowbreak{}\hskip.1em \RB{26.492}{7.766}{1.995}{%
\RG{0.000}{0.000}{10.9091}{mi}{110}\RG{9.579}{0.000}{10.9091}{sy}{21}\RG{21.094}{0.000}{10.9091}{rm}{50}%
}\hskip.1em\allowbreak{} 。当\allowbreak{}\hskip.1em \RB{26.148}{7.766}{0.620}{%
\RG{0.000}{0.000}{10.9091}{mi}{110}\RG{9.579}{0.000}{10.9091}{rm}{61}\RG{21.077}{0.000}{10.9091}{rm}{49}%
}\hskip.1em\allowbreak{} 时虽有\allowbreak{}\hskip.1em \RB{28.625}{8.191}{0.740}{%
\RG{0.000}{0.000}{10.9091}{mi}{67}\RG{11.608}{0.000}{10.9091}{rm}{61}\RG{23.106}{0.000}{10.9091}{rm}{48}%
}\hskip.1em\allowbreak{} ，但\allowbreak{}\hskip.1em \RB{26.070}{8.071}{0.740}{%
\RG{0.000}{0.000}{10.9091}{mi}{107}\RG{9.053}{0.000}{10.9091}{rm}{61}\RG{20.551}{0.000}{10.9091}{rm}{48}%
}\hskip.1em\allowbreak{} 且\allowbreak{}\hskip.1em \RB{57.430}{9.759}{2.856}{%
\RG{0.000}{0.000}{10.9091}{mi}{67}\RG{8.577}{3.959}{7.9701}{rm}{48}\RG{16.340}{0.000}{10.9091}{rm}{61}\RG{27.855}{0.000}{10.9091}{mi}{73}\RG{32.650}{-1.636}{7.9701}{rm}{49}\RG{40.413}{0.000}{10.9091}{sy}{54}\RG{40.413}{0.000}{10.9091}{rm}{61}\RG{51.912}{0.000}{10.9091}{rm}{48}%
}\hskip.1em\allowbreak{} ；因此该问必须补充\allowbreak{}\hskip.1em \RB{26.492}{7.761}{2.000}{%
\RG{0.000}{-0.005}{10.9091}{mi}{110}\RG{9.578}{-0.005}{10.9091}{sy}{21}\RG{21.094}{-0.005}{10.9091}{rm}{50}%
}\hskip.1em\allowbreak{} 。

小结：交换子与极小多项式给出次数界；单一最大Jordan块的中心化子由矩阵多项式构成，给出\allowbreak{}\hskip.1em \RB{24.952}{7.764}{0.622}{%
\RG{0.000}{-0.002}{10.9091}{mi}{110}\RG{8.972}{-0.002}{10.9091}{sy}{0}\RG{19.881}{-0.002}{10.9091}{rm}{49}%
}\hskip.1em\allowbreak{} 界。

\[
\RB{278.451}{7.723}{2.297}{%
\RT{0.000}{-0.004}{7.9701}{题}\RT{8.966}{-0.004}{7.9701}{目}\RT{17.933}{-0.004}{7.9701}{来}\RT{26.899}{-0.004}{7.9701}{源}\RT{35.866}{-0.004}{7.9701}{：}\RT{44.832}{-0.004}{7.9701}{谢}\RT{53.799}{-0.004}{7.9701}{启}\RT{62.765}{-0.004}{7.9701}{鸿}\RT{71.732}{-0.004}{7.9701}{《}\RT{80.698}{-0.004}{7.9701}{复}\RT{89.665}{-0.004}{7.9701}{旦}\RT{98.631}{-0.004}{7.9701}{大}\RT{107.598}{-0.004}{7.9701}{学}\RT{116.564}{-0.004}{7.9701}{高}\RT{125.531}{-0.004}{7.9701}{等}\RT{134.497}{-0.004}{7.9701}{代}\RT{143.464}{-0.004}{7.9701}{数}\RT{152.430}{-0.004}{7.9701}{习}\RT{161.397}{-0.004}{7.9701}{题}\RT{170.363}{-0.004}{7.9701}{集}\RT{179.329}{-0.004}{7.9701}{（}\RT{188.296}{-0.004}{7.9701}{2}\RT{192.779}{-0.004}{7.9701}{0}\RT{197.262}{-0.004}{7.9701}{2}\RT{201.745}{-0.004}{7.9701}{3}\RT{209.214}{-0.004}{7.9701}{版}\RT{218.180}{-0.004}{7.9701}{）}\RT{222.672}{-0.004}{7.9701}{》}\RT{227.156}{-0.004}{7.9701}{，}\RT{236.122}{-0.004}{7.9701}{第}\RT{248.066}{-0.004}{7.9701}{7}\RT{252.549}{-0.004}{7.9701}{5}\RT{260.018}{-0.004}{7.9701}{页}\RT{268.985}{-0.004}{7.9701}{。}%
}
\]

\end{document}
